\documentclass[11pt]{amsart}

\usepackage{amsmath}
\usepackage{amssymb}
\usepackage{amsthm}
\usepackage{booktabs}
\usepackage{array}
\usepackage{hyperref}

% ---------------------------------------------------------------------------
% Theorem environments
% ---------------------------------------------------------------------------
\newtheorem{theorem}{Theorem}[section]
\newtheorem{proposition}[theorem]{Proposition}
\newtheorem{lemma}[theorem]{Lemma}
\newtheorem{corollary}[theorem]{Corollary}
\theoremstyle{definition}
\newtheorem{definition}[theorem]{Definition}
\theoremstyle{remark}
\newtheorem{remark}[theorem]{Remark}

\numberwithin{equation}{section}
\setlength{\emergencystretch}{2em}

% ---------------------------------------------------------------------------
% Macros
% ---------------------------------------------------------------------------
\newcommand{\NN}{\mathbb{N}}
\newcommand{\ZZ}{\mathbb{Z}}
\newcommand{\QQ}{\mathbb{Q}}
\newcommand{\RR}{\mathbb{R}}
\newcommand{\CC}{\mathbb{C}}
\newcommand{\FF}{\mathbb{F}}
\newcommand{\cA}{\mathcal{A}}
\newcommand{\cR}{\mathcal{R}}
\newcommand{\eps}{\varepsilon}
\newcommand{\ind}[1]{[\,#1\,]}
\newcommand{\lev}{\operatorname{lev}}
\newcommand{\cb}{\operatorname{cb}}
\newcommand{\red}{\operatorname{red}}
\newcommand{\Li}{\operatorname{Li}}
\newcommand{\Ed}{\deg_E}
\newcommand{\lean}[1]{{\def\_{\textunderscore\allowbreak}\texttt{#1}}}

\begin{document}

\title{\texorpdfstring{$\chi_2(1/3)$ and $\chi_2(1/5)$ are irrational}{chi2(1/3) and chi2(1/5) are irrational}}

\author{Shawn Krisman}

\subjclass[2020]{Primary 11J72; Secondary 33B30, 41A21, 31A15, 68V20}
\keywords{Legendre chi function, dilogarithm, irrationality, Hankel determinants, Pad\'e approximation,
logarithmic potential theory, formal verification, Lean}

\date{\today}

\begin{abstract}
Let $\chi_2(z)=\sum_{n\ge0}z^{2n+1}/(2n+1)^2$ be Legendre's chi function. We prove that $\chi_2(1/3)$ and
$\chi_2(1/5)$ are irrational. Equivalently, $\Li_2(1/q)-\Li_2(-1/q)\notin\QQ$ for $q=3,5$. The proof uses
the Hankel determinants $\Delta_h(\chi_2(1/q))$ of the moment functional
$g\mapsto\sum_{t\ge1}q^{-2t}\bigl(g(u)\,u/D_h(u)\bigr)'|_{u=t}$, where $D_h(u)=\prod_{k<h}(u+k+\frac12)$.
These determinants are values of polynomials $\Delta_h(X)\in\QQ[X]$ of degree at most $h$. We bound the
denominators of $\Delta_h(X)$ prime by prime: a bound adapted to residue classes for the large primes, a
Lagrange-basis bound for the small ones, and separate arguments for the primes $2$, $3$ and $q$. We bound
the size of $\Delta_h(\chi_2(1/q))$ by a discrete log-gas argument against an explicit dual certificate,
and prove that $\Delta_h(a/b)\ne0$ for every rational $a/b$ and infinitely many $h$, through a reduction modulo
an auxiliary prime $\ell=h/3$. The margins are below $-0.054$ ($q=3$) and $-0.256$ ($q=5$) per $h^2$. The whole
proof, including the prime number theorem and both numerical certificates, has been formalised in Lean~4 with Mathlib
and checked by the Lean kernel, using only Lean's three standard axioms.
\end{abstract}

\maketitle

\tableofcontents

% ===========================================================================
\section{Introduction}\label{sec:intro}
% ===========================================================================

For $|z|\le1$ let
\[
  \chi_2(z)=\sum_{n=0}^{\infty}\frac{z^{2n+1}}{(2n+1)^2}=\tfrac12\bigl(\Li_2(z)-\Li_2(-z)\bigr)
\]
be Legendre's chi function \cite[\S1.8]{Lewin1981}.

\begin{theorem}\label{thm:main}
The numbers $\chi_2(1/3)=0.33762317823243597996\ldots$ and $\chi_2(1/5)=0.20090195660862019582\ldots$ are
irrational.
\end{theorem}

\subsection{Context}
We know of no earlier proof that $\chi_2(1/q)$ is irrational for any small integer $q\ge2$.
\begin{itemize}
\item \emph{Large $q$.} We have $\chi_2(1/q)=(4q)^{-1}\Phi(q^{-2},2,\tfrac12)=\frac q4L(\tfrac12,2;q^{-2})$, where
  $\Phi$ is Lerch's transcendent and $L(\alpha,k;x)=\sum_{n\ge1}x^n/(n-\alpha)^k$. Hata \cite{Hata1990} proved
  linear independence measures for values $L(\alpha_j,k_j;1/Q)$ with rational shifts, for $|Q|>Q_0$ with explicit $Q_0$.
  The application to $\chi_2(1/q)$ has $Q=q^2$, so $Q=9$ and $Q=25$ here. Hata's headline application,
  $\Li_2(1/Q)\notin\QQ$ for $Q\ge12$ or $Q\le-8$, is for the shift $0$, which is easier than the shift $\frac12$,
  so $q=3$ ($Q=9$) is presumably outside his range. We have not determined his $Q_0$ for the shift $\frac12$, which
  decides whether his method covers $q=5$ ($Q=25$).
  Heuristically the half shift costs a factor of about $e^2$ in $Q$, because the denominators become
  $\operatorname{lcm}$ of the odd numbers up to $2n$, which suggests that it is not covered [heuristic]. The explicit
  criterion of David, Hirata-Kohno and Kawashima \cite[Thm.~2.1]{DHK2022} for the shift $\frac12$, depths $\le2$, at the
  single point $1/q^2$ requires $\log q^2>2\log2+6\log3+8$, i.e.\ $q\ge2949$. Their multi-point criterion
  \cite[Thm.~2.1]{DHK2020} does not apply to $\chi_2(z)=\Li_2(z)-\frac14\Li_2(z^2)$ at the two points $1/q$, $1/q^2$ for
  any $q$, because it requires all points to be of comparable size.
\item \emph{Single dilogarithm values do not suffice.} Hata \cite{Hata1993} proved $\Li_2(-1/5)\notin\QQ$. The
  irrationality of $\Li_2(1/5)$ has been claimed recently by Prellberg \cite{Prellberg2026} and in S.~Liu's Lean
  repository \cite{LiuLi2}; neither is refereed. In any case two irrational numbers can have a rational
  difference, so neither result says anything about $\chi_2(1/5)=\frac12(\Li_2(\frac15)-\Li_2(-\frac15))$. The
  linear independence results of Rhin--Viola \cite{RhinViola2005} (one point $1/z$) and Viola--Zudilin
  \cite{ViolaZudilin2018} (the two points $1/z$ and $1/(1-z)$) do not cover the pair $\pm\frac15$ either.
\item \emph{Identities do not suffice.} Landen's identity \cite[App.~A.2.1, (24)]{Lewin1981}
  $\chi_2\bigl(\tfrac{1-x}{1+x}\bigr)+\chi_2(x)=\tfrac{\pi^2}{8}+\tfrac12\log x\,\log\tfrac{1+x}{1-x}$ gives, for
  example, $\chi_2(\frac13)+\chi_2(\frac12)=\frac{\pi^2}8-\frac12\log2\log3$. Both $\chi_2(\frac12)$ and $\log2\log3$
  are open.
\end{itemize}
Our search for prior work covered the web, arXiv and recent Lean repositories, but not MathSciNet or zbMATH.

\subsection{The method}\label{sec:method}
Fix an integer $q\ge2$ and put $w:=q^{-2}$. For $h\ge1$ let $x_k:=-(k+\frac12)$ for $0\le k<h$,
$D(u):=\prod_{k<h}(u-x_k)$ and
\begin{equation}\label{eq:Lambda}
  \Lambda(g):=\sum_{t\ge1}w^t\Bigl(\frac{g(u)\,u}{D(u)}\Bigr)'\Big|_{u=t}\qquad(g\in\QQ[u]).
\end{equation}
Partial fractions show that $\Lambda(g)=A(g)+\chi_2(1/q)\,B(g)$ with explicit rational functionals $A$ and $B$
(Lemma~\ref{lem:pole}). Hence the Hankel determinant $\det(\Lambda(u^{i+j}))_{i,j<h}$ is the value at
$X=\chi_2(1/q)$ of
\[
  \Delta_h(X):=\det\bigl(A(u^{i+j})+X\,B(u^{i+j})\bigr)_{0\le i,j<h}\in\QQ[X],\qquad \deg\Delta_h\le h.
\]
This is a Pad\'e-type construction of the kind that Q.~Tang used for $\Li_2(1/2)$ \cite{TangLi2}; the
half-integer poles are what make $\chi_2$ appear. The proof has three inputs.
\begin{enumerate}
\item \emph{Denominators} (Sections~\ref{sec:arith1}--\ref{sec:arith2}). There are rationals $d_h>0$ with
  $d_h\Delta_h\in\ZZ[X]$ and
  \[
    \limsup_{h\to\infty}\frac{\log d_h}{h^2}\le\alpha_q:=\frac{133}{60}+\frac{\log2}{4}+\ind{q=5}\frac{\log3}{4}-\frac{\log q}{3}.
  \]
  The term $133/60=\frac45+\frac{17}{12}$ comes from the primes $p\nmid 2q(q^2-1)$. Primes $p<2h/5$ are handled
  by a Lagrange-basis bound and primes $p\ge2h/5$ by a bound adapted to the residue classes modulo $p$; the
  prime number theorem then sums the two contributions. The prime $q$ gives a \emph{gain}: a reflection of the functional
  shows that every coefficient of $\Delta_h$ has $v_q\ge\lfloor(h^2-7h)/3\rfloor-O(h\log h)$. The primes $2$ and
  $3\mid q^2-1$ occur in the denominators with exponents at most $\frac14h^2+O(h\log h)$.
\item \emph{Size} (Sections~\ref{sec:analytic}--\ref{sec:cert}). For every \emph{dual certificate} (a step
  density $\mu$ with cell bounds on its potential) $\limsup h^{-2}\log|\Delta_h(\chi_2(1/q))|\le\Phi$, an
  explicit exponent computed from the certificate.
  The proof treats $\Delta_h(\chi_2(1/q))$ as the partition function of a discrete log-gas with at most two
  particles per integer node, and compares it with $\mu$ through a lattice energy inequality of P\'olya type.
  Explicit certificates, checked by exact rational interval arithmetic, give $\Phi<-2.078$ for $q=3$ and
  $\Phi<-2.385$ for $q=5$.
\item \emph{Nonvanishing} (Section~\ref{sec:nonvan}). For every rational $a/b$, $\Delta_h(a/b)\ne0$ for infinitely
  many $h$. We take $h=3\ell$ with $\ell$ a large prime. In a Lagrange basis and after an $\ell$-adic rescaling,
  the Gram matrix of $\Lambda_{a/b}=A+\frac ab B$ is $\ell$-integral and block upper triangular modulo~$\ell$, with
  explicit $3\times3$ diagonal blocks of three types. Their determinants are three fixed nonzero rationals
  for each $q$, so they are units for all large $\ell$.
\end{enumerate}
Since $\alpha_3=2.0237\ldots$ and $\alpha_5=2.1281\ldots$, the margins are
\[
  \alpha_3-2.078<-0.054,\qquad \alpha_5-2.385<-0.256,
\]
and the standard argument (Proposition~\ref{prop:criterion}) concludes. Table~\ref{tab:budget} lists the
contributions.

\begin{table}[ht]
\centering
\begin{tabular}{@{}lccl@{}}
\toprule
contribution (per $h^2$) & $q=3$ & $q=5$ & source\\
\midrule
primes $p\nmid2q(q^2-1)$: $\frac{133}{60}$ & $2.21667$ & $2.21667$ & Thm.~\ref{thm:generic}\\
$p=2$: $\frac14\log2$ & $0.17329$ & $0.17329$ & Lemma~\ref{lem:two}\\
$p=3\mid q^2-1$: $\frac14\log3$ & --- & $0.27465$ & Lemma~\ref{lem:poleorder}\\
$p=q$ (gain): $-\frac13\log q$ & $-0.36620$ & $-0.53648$ & Lemma~\ref{lem:qadic}\\
\midrule
arithmetic $\alpha_q$ & $2.02375$ & $2.12813$ & \\
analytic $\Phi$ & $<-2.07832$ & $<-2.38515$ & Thms.~\ref{thm:B}, \ref{thm:cert}\\
\midrule
margin $\alpha_q+\Phi$ & $<-0.0545$ & $<-0.2570$ & \\
\bottomrule
\end{tabular}
\caption{The exponential budget. The analytic values are those of the certificates of
Section~\ref{sec:cert}; the main theorem only uses $\Phi<-2.078$ and $\Phi<-2.385$.}\label{tab:budget}
\end{table}

\subsection{The formalisation}
The statement of Theorem~\ref{thm:main} and its complete proof are formalised in Lean~4 \cite{LeanFour} with
Mathlib \cite{Mathlib} and checked by the Lean kernel. This includes the prime number theorem, the two dual
certificates and the block determinants. The proof depends only on Lean's three standard axioms
(\lean{propext}, \lean{Classical.choice}, \lean{Quot.sound}). It contains no \lean{sorry} and no
\lean{native\_decide}: every proof, including every computation, is checked by the kernel. The repository is
\url{https://github.com/telaviv/chi2-irrationality} \cite{Chi2Repo}; Section~\ref{sec:lean} describes it.

The prose proof below is self-contained and can be read without the formalisation, except that two
kinds of finite computation are only described, not reproduced: the two dual certificates (Theorem~\ref{thm:cert}),
which are decidable statements about rational data proved by \lean{decide +kernel}, and the six block
determinants of Section~\ref{sec:nonvan}, which are evaluated exactly by \lean{norm\_num}.

\subsection{Organisation and notation}
Section~\ref{sec:construction} sets up the functional and proves the pole expansion. Section~\ref{sec:criterion}
states the criterion and reduces Theorem~\ref{thm:main} to the three inputs.
Sections~\ref{sec:arith1} and~\ref{sec:arith2} bound the denominators, Sections~\ref{sec:analytic}
and~\ref{sec:cert} bound the size, and Section~\ref{sec:nonvan} proves nonvanishing.
Section~\ref{sec:lean} describes the formalisation and Section~\ref{sec:remarks} contains remarks, including
why $\chi_2(1/2)$ is out of reach of this construction.

Throughout, $p$ and $\ell$ denote primes; $v_p$ is the $p$-adic valuation on $\QQ$ and $|x|_p=p^{-v_p(x)}$.
$\lfloor\cdot\rfloor$ and $\lceil\cdot\rceil$ are the floor and ceiling, $(x)_+=\max(x,0)$, and $\ind{P}$ is $1$ if
$P$ holds and $0$ otherwise. For a sequence $f$ on $\ZZ_{\ge0}$, $\Delta f(t)=f(t+1)-f(t)$ is the forward
difference. We write $\theta(x)=\sum_{p\le x}\log p$.

% ===========================================================================
\section{The construction}\label{sec:construction}
% ===========================================================================

\subsection{Data}\label{sec:data}
Fix an integer $q\ge2$; later $q\in\{3,5\}$. Put
\[
  z:=1/q,\qquad w:=z^2=q^{-2},\qquad \chi:=\chi_2(z).
\]
For $h\ge1$ we use $h$ poles $x_k:=-(k+\frac12)$, $0\le k<h$, and put
\[
  D(u):=\prod_{k<h}(u-x_k),\qquad \omega_k:=\frac{x_k}{D'(x_k)},
\]
where $D'(x_k)=\prod_{m\ne k}(x_k-x_m)=\prod_{m\ne k}(m-k)=(-1)^kk!\,(h-1-k)!$. We also need the partial sums
\[
  S_k(q):=\sum_{j=0}^{k}\frac{q^{2(k-j)}}{(2j+1)^2}=q^{2k+1}T(2k+1),\qquad T(M):=\sum_{\substack{s\le M\\ s\text{ odd}}}\frac{z^s}{s^2},
\]
with $T(M)=0$ for $M<1$. For a rational function $F$ regular at the positive integers and of polynomial growth put
\[
  \mu(F):=\sum_{t\ge1}w^tF'(t),
\]
an absolutely convergent series since $0<w<1$. The functional \eqref{eq:Lambda} is $\Lambda(g)=\mu(g\varphi)$ with
$\varphi(u):=u/D(u)$.

\subsection{The polynomial part}
For a polynomial $R$, $\mu(R)$ is rational. Indeed $\mu(u^d)=d\sum_{t\ge1}t^{d-1}w^t$, and the power sums
$\sigma_n:=\sum_{t\ge1}t^nw^t$ satisfy $\sigma_0=w/(1-w)$ and
\begin{equation}\label{eq:powsum}
  (1-w)\,\sigma_{n+1}=\sum_{j=0}^{n}\binom{n+1}{j}(-1)^{n-j}\sigma_j ,
\end{equation}
by expanding $\sum_t t^{n+1}w^t-\sum_t(t-1)^{n+1}w^{t}=(1-w)\sigma_{n+1}$. We shall also use the Newton form
\begin{equation}\label{eq:newton}
  \mu(R)=\sum_{b\ge0}\Delta^b[R'](0)\,\Omega_b,\qquad
  \Omega_b:=\sum_{t\ge1}\binom tbw^t=\begin{cases}\dfrac{1}{q^2-1},&b=0,\\[6pt] \dfrac{q^2}{(q^2-1)^{b+1}},&b\ge1,\end{cases}
\end{equation}
which follows from $R'(t)=\sum_b\Delta^b[R'](0)\binom tb$. Here
\[
  \Delta^b[R'](0)=\sum_{k\le b}(-1)^{b-k}\binom bk\,R'(k),
\]
and the sum over $b$ is finite ($b<\deg R$).

\subsection{The pole expansion}

\begin{lemma}[Lemma~1, pole expansion]\label{lem:pole}
For $g\in\QQ[u]$ write $g(u)u=D(u)R_g(u)+r_g(u)$ with $\deg r_g<h$. Then
\[
  \Lambda(g)=A(g)+\chi\,B(g),
\]
where
\[
  A(g):=\mu(R_g)+4\sum_{k<h}g(x_k)\,\omega_k\,S_k(q),\qquad
  B(g):=-4\sum_{k<h}g(x_k)\,\omega_k\,q^{2k+1}.
\]
Both are rational and $\QQ$-linear in $g$.
\end{lemma}

\begin{proof}
By Lagrange interpolation at the $h$ nodes $x_k$,
\[
  r_g(u)=\sum_k\frac{r_g(x_k)\,D(u)}{D'(x_k)\,(u-x_k)},\qquad r_g(x_k)=g(x_k)\,x_k .
\]
Hence $g\varphi=R_g+\sum_kg(x_k)\omega_k/(u-x_k)$. For the pole terms,
\[
  \mu\Bigl(\frac1{u-x_k}\Bigr)=-\sum_{t\ge1}\frac{w^t}{(t+k+\frac12)^2}=-4\sum_{t\ge1}\frac{z^{2t}}{(2t+2k+1)^2}.
\]
With $s=2t+2k+1$ (odd, $s\ge2k+3$) and $z^{2t}=q^{2k+1}z^s$ this equals
$-4q^{2k+1}\bigl(\chi-T(2k+1)\bigr)=-4\bigl(q^{2k+1}\chi-S_k(q)\bigr)$. All series converge absolutely, so we may
sum termwise.
\end{proof}

For an indeterminate $X$ put $\Lambda_X:=A+X\,B$ and
\[
  H(X):=\bigl(\Lambda_X(u^{i+j})\bigr)_{0\le i,j<h},\qquad \Delta_h(X):=\det H(X)\in\QQ[X].
\]
The entries are linear in $X$, so $\deg\Delta_h\le h$, and by Lemma~\ref{lem:pole}
\begin{equation}\label{eq:value}
  \Delta_h(\chi)=\det\bigl(\Lambda(u^{i+j})\bigr)_{0\le i,j<h}.
\end{equation}

\subsection{Change of basis}
If $a_0,\dots,a_{h-1}$ and $b_0,\dots,b_{h-1}$ are polynomials of degree $<h$ with coefficient matrices $C_a$ and $C_b$
(rows indexed by $i$, columns by the monomials $u^0,\dots,u^{h-1}$), then by linearity of $\Lambda_X$
\begin{equation}\label{eq:gram}
  \bigl(\Lambda_X(a_ib_j)\bigr)_{i,j}=C_a\,H(X)\,C_b^{\mathsf T},\qquad
  \det\bigl(\Lambda_X(a_ib_j)\bigr)=\det C_a\det C_b\,\Delta_h(X).
\end{equation}
We use three bases. The \emph{Lagrange basis} $f_i:=D/(u-x_i)$ has $f_i(x_k)=\ind{i=k}D'(x_i)$. Write $V=(x_k^d)_{d,k<h}$
for the Vandermonde matrix. Then $C_fV=\operatorname{diag}(D'(x_i))$, and
$\prod_iD'(x_i)=\pm(\det V)^2$ with $\det V=\prod_{k<k'}(x_{k'}-x_k)=\pm\prod_{k<k'}(k'-k)$. The \emph{anchored basis}
\[
  a_i:=\prod_{m<i}(u-x_m),\qquad b_j:=\prod_{h-j\le m<h}(u-x_m)
\]
consists of monic polynomials of degrees $i$ and $j$, so $C_a$ and $C_b$ are unitriangular and
$\det(\Lambda_X(a_ib_j))=\Delta_h(X)$. The third basis, adapted to residue classes, is introduced in
Section~\ref{sec:classbound}.

% ===========================================================================
\section{The criterion and the structure of the proof}\label{sec:criterion}
% ===========================================================================

\begin{proposition}[Lemma~E]\label{prop:criterion}
Let $\xi\in\RR$, and for $h\ge1$ let $P_h\in\QQ[X]$ with $\deg P_h\le h$ and let $d_h>0$ be rationals with
$d_hP_h\in\ZZ[X]$. Suppose that there is $c>0$ with $d_h|P_h(\xi)|\le e^{-ch^2}$ for all large $h$, and that for
every $a\in\ZZ$ and $b\ge1$ we have $P_h(a/b)\ne0$ for infinitely many $h$. Then $\xi$ is irrational.
\end{proposition}

\begin{proof}
Suppose $\xi=a/b$. For every $h$, $N_h:=b^hd_hP_h(a/b)$ is an integer, because $d_hP_h$ has integer coefficients
and degree $\le h$. Choose $h$ with $P_h(a/b)\ne0$, $d_h|P_h(\xi)|\le e^{-ch^2}$ and $b^he^{-ch^2}<1$. Then
$1\le|N_h|<1$, a contradiction.
\end{proof}

\begin{theorem}\label{thm:conditional}
Let $q\ge2$. Suppose there are rationals $d_h>0$ and real numbers $\alpha,\beta$ with $\alpha+\beta<0$ such that
\begin{enumerate}
\item $d_h\Delta_h\in\ZZ[X]$ for every $h$, and for every $\eps>0$, $\log d_h\le(\alpha+\eps)h^2$ for all large $h$;
\item for every $\eps>0$, $|\Delta_h(\chi_2(1/q))|\le e^{(\beta+\eps)h^2}$ for all large $h$;
\item for every $a\in\ZZ$ and $b\ge1$, $\Delta_h(a/b)\ne0$ for infinitely many $h$.
\end{enumerate}
Then $\chi_2(1/q)$ is irrational.
\end{theorem}

\begin{proof}
Apply Proposition~\ref{prop:criterion} with $P_h=\Delta_h$ and $c=-(\alpha+\beta)/2$, using (i) and (ii) with
$\eps=c/2$.
\end{proof}

For $q=3$ and $q=5$ the three hypotheses are Theorem~\ref{thm:arith} with $\alpha=\alpha_q$,
Theorems~\ref{thm:B} and~\ref{thm:cert} with $\beta=-2.078$ and $\beta=-2.385$, and
Theorem~\ref{thm:nonvan}. The final assembly is in Section~\ref{sec:proof}.

% ===========================================================================
\section{Denominators I: the symbolic machine and the generic primes}\label{sec:arith1}
% ===========================================================================

\subsection{The free ring}\label{sec:free}
It is convenient to treat $q$ and $1/(q^2-1)$ as independent indeterminates. Let $\cR:=\QQ[q,E]$ be the polynomial
ring in two variables. In $\cR$ put $\Omega_0:=E$ and $\Omega_b:=q^2E^{b+1}$ for $b\ge1$, and for $R\in\QQ[u]$ define
$\mu^{\cR}(R):=\sum_b\Delta^b[R'](0)\,\Omega_b\in\cR$ as in \eqref{eq:newton}. Replacing $\mu$ by $\mu^{\cR}$, and
reading $S_k(q)$ and $q^{2k+1}$ as elements of $\QQ[q]$, Lemma~\ref{lem:pole} defines functionals $A^{\cR},B^{\cR}$ and
\[
  \Delta^{\cR}_h(X):=\det\bigl(A^{\cR}(u^{i+j})+X\,B^{\cR}(u^{i+j})\bigr)_{i,j<h}\in\cR[X].
\]
For an integer $q_0\ge2$, the ring homomorphism $\operatorname{ev}_{q_0}\colon\cR\to\QQ$ with $q\mapsto q_0$ and
$E\mapsto1/(q_0^2-1)$ maps $\Omega_b$ to the value \eqref{eq:newton}, hence $\mu^{\cR}$ to $\mu$ and
$\Delta^{\cR}_h$ to $\Delta_h$ (computed with $q=q_0$).

\emph{Gauss norms.} For a prime $p$ and $f\in\cR$ let $\|f\|_p$ be the maximum of $|c|_p$ over the rational
coefficients $c$ of $f$, and extend this coefficientwise to $\cR[X]$. Then $\|f+g\|_p\le\max(\|f\|_p,\|g\|_p)$ and
$\|fg\|_p\le\|f\|_p\|g\|_p$. If $|v_0|_p\le1$ and $|v_1|_p\le1$, the evaluation $q\mapsto v_0$, $E\mapsto v_1$
does not increase the norm. By the Leibniz expansion:

\begin{lemma}\label{lem:leibniz}
Let $M$ be an $h\times h$ matrix over $\cR[X]$ with $\|M_{ij}\|_p\le\rho_i\kappa_j$ for some $\rho_i,\kappa_j\ge0$. Then
$\|\det M\|_p\le\prod_i\rho_i\prod_j\kappa_j$.
\end{lemma}

Two facts carry all the dependence on $q$:
\begin{itemize}
\item[(F1)] $\|\Omega_b\|_p=1$ for all $b$ and $p$, so for $R\in\QQ[u]$ we have
  $\|\mu^{\cR}(R)\|_p\le\sup_{t\in\ZZ_{\ge0}}|R'(t)|_p$, because $\Delta^b[R'](0)$ is an integer combination of
  $R'(0),\dots,R'(b)$;
\item[(F2)] $\|q^{2k+1}X-S_k(q)\|_p\le p^{2\#\{e\ge1:\,p^e\le2k+1\}}$, since the coefficients of $S_k$ are $1/s^2$
  with $s$ odd, $s\le2k+1$, and $v_p(s)\le\#\{e\ge1:p^e\le s\}$.
\end{itemize}

\emph{The Gram identity.} Write $\Lambda^{\cR}_X:=A^{\cR}+XB^{\cR}$, a $\QQ$-linear map $\QQ[u]\to\cR[X]$. By the argument
of \eqref{eq:gram}, for any family $(e_m)_{m<h}$ of polynomials of degree $<h$ with coefficient matrix $C_e$,
\begin{equation}\label{eq:gramR}
  \det\bigl(\Lambda^{\cR}_X(e_me_{m'})\bigr)_{m,m'<h}=(\det C_e)^2\,\Delta^{\cR}_h(X).
\end{equation}
For a pole $x_l$ we have, by the proof of Lemma~\ref{lem:pole}, $\Lambda^{\cR}_X(g)=\mu^{\cR}(R)-4c\,(q^{2l+1}X-S_l(q))$
whenever $g\varphi=R+c/(u-x_l)$ with $R$ a polynomial; more generally the pole parts add up.

Below, $\lev(p):=\#\{e\ge1:p^e\le2h-1\}$. For $t\in\ZZ_{\ge0}$ and $m<h$ we have $t-x_m=(2t+2m+1)/2$, and for $p$ odd
$v_p(t-x_m)=v_p(2t+2m+1)$. The numbers $2t+2m+1$, $m<h$, are $h$ consecutive odd integers, so for each level
$d=p^e$ those divisible by $d$ are the $m$ in one residue class modulo $d$.

\subsection{All odd primes: the Lagrange basis}

\begin{lemma}[Lemma~9]\label{lem:nine}
For every odd prime $p$ and every $h\ge1$, $\|\Delta^{\cR}_h\|_p\le p^{2h\lev(p)}$.
\end{lemma}

\begin{proof}
Use the Lagrange basis $f_i=D/(u-x_i)$ and put $G_{ij}:=\Lambda^{\cR}_X(f_if_j)$. Since $C_fV=\operatorname{diag}(D'(x_i))$
and $\prod_iD'(x_i)=\pm(\det V)^2$, we get $(\det C_f)^2=\pm\prod_iD'(x_i)$, so by \eqref{eq:gramR}
\[
  \|\Delta^{\cR}_h\|_p=\|\det G\|_p\Big/\prod_i|D'(x_i)|_p .
\]
By Lemma~\ref{lem:leibniz} it suffices to prove the \emph{row claim} $\|G_{ij}\|_p\le|D'(x_i)|_p\,p^{2\lev(p)}$. For a
level $e\ge1$ put $N_e(i):=\#\{m<h:m\ne i,\ p^e\mid m-i\}$; then $v_p(D'(x_i))=\sum_eN_e(i)$, and $N_e(i)=0$ unless
$p^e\le h-1$. So it suffices to show that every value $R'(t)$ entering (F1) has
\begin{equation}\label{eq:rowclaim}
  v_p(R'(t))\ge\sum_e\bigl(N_e(i)-2\ind{N_e(i)\ge1}\bigr),
\end{equation}
and that the pole part obeys the same bound.

Fix $t\in\ZZ_{\ge0}$ and a level $d=p^e$, and let $m_0$ be the residue class mod $d$ of the indices $m$ with
$d\mid2t+2m+1$. The class of $i$ has $N_e(i)+1$ elements in $[0,h)$, and residue classes mod $d$ in $[0,h)$ differ in
size by at most one, so the class $m_0$ has at least $N_e(i)$ elements.

\emph{Off-diagonal entries} ($i\ne j$). Here $f_if_j\varphi=F_{ij}:=u\prod_{m\ne i,j}(u-x_m)$ is a polynomial, so
$G_{ij}=\mu^{\cR}(F_{ij})$. Each term of $F'_{ij}(t)$ is a product of all factors but one. At level $d$ the factors
$t-x_m$ divisible by $d$ are those with $m\equiv m_0$, $m\ne i,j$; omitting one factor leaves at least
$N_e(i)-2$ of them in both cases $m_0\equiv i$ and $m_0\not\equiv i$. Summing over levels ($v_p$ of a product is
$\sum_e$ of the number of factors divisible by $p^e$) gives \eqref{eq:rowclaim}.

\emph{Diagonal entries.} Put $g_i:=u\prod_{m\ne i}(u-x_m)$. Then $f_i^2\varphi=g_i/(u-x_i)=Q_i+g_i(x_i)/(u-x_i)$ with
$Q_i:=(g_i(u)-g_i(x_i))/(u-x_i)$ and $g_i(x_i)=x_iD'(x_i)$. The pole part is $-4x_iD'(x_i)(q^{2i+1}X-S_i)$; by (F2),
$|x_i|_p\le1$ and $2i+1\le2h-1$, its norm is at most $|D'(x_i)|_p\,p^{2\lev(p)}$. For the polynomial part, list the
$h$ roots $\{x_m:m\ne i\}\cup\{0\}$ of $g_i$ as $\alpha_1,\dots,\alpha_h$; telescoping
$(u-\alpha_r)-(x_i-\alpha_r)=u-x_i$ gives
\[
  Q_i(u)=\sum_{r=1}^{h}\prod_{s<r}(x_i-\alpha_s)\prod_{s>r}(u-\alpha_s).
\]
Order the roots by decreasing $v_p(m-i)$, with $0$ last. Then for every level $e$ simultaneously, the $\alpha_s=x_m$ with
$p^e\mid m-i$ form an initial segment of length $N_e(i)$ (note $x_i-x_m=m-i$). Fix $t$, a level $d=p^e$ and a term
$r$ of $Q_i'(t)$ (a product with one more factor omitted). If $r>N_e(i)$, the first product already has $N_e(i)$
factors divisible by $d$. If $r\le N_e(i)$, it has $r-1$; if moreover $m_0\equiv i$, the remaining $N_e(i)-r$ class
members sit at positions $s>r$ and at most one of them is omitted, giving at least $N_e(i)-2$ in total; if
$m_0\not\equiv i$, the whole class $m_0$ (at least $N_e(i)$ members, all with $v_p(m-i)<e$) sits at positions
$s>N_e(i)\ge r$, giving at least $N_e(i)-1$. So \eqref{eq:rowclaim} holds for $Q_i'(t)$ as well. Finally
$\#\{e:N_e(i)\ge1\}\le\lev(p)$.
\end{proof}

\subsection{One-level primes: a basis adapted to residue classes}\label{sec:classbound}
For a modulus $d$ and $m<h$ write $r_m:=\lfloor m/d\rfloor$ (the rank of $m$ in its class), $n_m$ for the size of the
class of $m$ modulo $d$ in $[0,h)$, and put
\begin{gather*}
  w_d(l):=2\ind{d\le2l+1}-\ind{d\mid2l+1},\\
  M_m:=\max\{w_d(l):l\le m,\ l\equiv m\ (d)\},\\
  \cb(h,d):=\sum_{m<h}\bigl(2r_m+1+M_m-n_m\bigr)_+ .
\end{gather*}

\begin{proposition}[class bound]\label{prop:classbound}
Let $p$ be an odd prime with $p^2>2h$. Then $\|\Delta^{\cR}_h\|_p\le p^{\cb(h,p)}$.
\end{proposition}

\begin{proof}
Since $p^2>2h$, every difference $m-l$ of distinct indices and every odd $s\le2h-1$ has $v_p\le1$; so
$v_p(m-l)=\ind{p\mid m-l}$ and $v_p(2l+1)=\ind{p\mid2l+1}$. For $m<h$ let $\operatorname{seg}(m):=\{l\le m:l\equiv m\ (p)\}$
(it has $r_m+1$ elements), and put
\[
  e_m:=s_m\,\frac{D(u)}{\prod_{l\in\operatorname{seg}(m)}(u-x_l)},\qquad s_m:=\prod_{l\in\operatorname{seg}(m),\,l<m}(x_m-x_l).
\]
Then $\deg e_m=h-r_m-1<h$ and $v_p(s_m)=r_m$. The polynomial $e_m$ vanishes at $x_k$ unless $k\in\operatorname{seg}(m)$, and
$e_m(x_m)=D'(x_m)$. So the evaluation matrix $(e_m(x_k))=C_eV$ is lower triangular with diagonal $D'(x_m)$, hence
$\det C_e\det V=\prod_mD'(x_m)=\pm(\det V)^2$ and $\det C_e=\pm\det V$. By \eqref{eq:gramR},
$\|\Delta^{\cR}_h\|_p=\|\det G^e\|_p/|\det V|_p^2$ with $G^e_{mm'}:=\Lambda^{\cR}_X(e_me_{m'})$, and
$2v_p(\det V)=2\#\{k<k':p\mid k'-k\}=2\sum_mr_m$.

Let $\cA:=\operatorname{seg}(m)$, $\mathcal B:=\operatorname{seg}(m')$, $r=r_m$, $r'=r_{m'}$. Then
\[
  e_me_{m'}\varphi=s_ms_{m'}\,u\,\frac{\prod_{l\notin\cA\cup\mathcal B}(u-x_l)}{\prod_{l\in\cA\cap\mathcal B}(u-x_l)} .
\]
If $m$ and $m'$ lie in different classes, $\cA\cap\mathcal B=\emptyset$, the right side is $s_ms_{m'}$ times a polynomial in
$\ZZ[\frac12][u]$, and (F1) gives $v_p(G^e_{mm'})\ge r+r'$. If they lie in the same class, of size $n$, say with
$r\le r'$, then $\cA\subseteq\mathcal B$. Division by the monic polynomial $\prod_{l\in\cA}(u-x_l)\in\ZZ[\frac12][u]$ leaves a
$p$-integral polynomial part, again of valuation $\ge r+r'$. The residue at $x_l$, $l\in\cA$, is
\begin{gather*}
  c_l=s_ms_{m'}\,x_l\,\frac{\prod_{k\notin\mathcal B}(x_l-x_k)}{\prod_{k\in\cA,\,k\ne l}(x_l-x_k)},\\
  v_p(c_l)=r+r'+\ind{p\mid2l+1}+(n-1-r')-r,
\end{gather*}
because exactly the $n-1-r'$ class members after $m'$ lie outside $\mathcal B$, and the $r$ other elements of $\cA$ lie in
the class. By (F2) and the single level, the pole factor $q^{2l+1}X-S_l$ has norm at most $p^{2\ind{p\le2l+1}}$. Hence
\[
  v_p(G^e_{mm'})\ge r+r'+\beta(r,r'),\qquad \beta(a,b):=\min\bigl(0,\ n-1-a-b-\widehat M_{\min(a,b)}\bigr),
\]
where $\widehat M_a$ is the largest weight $w_p(l)$ among the class members of rank $\le a$. Since $\widehat M$ is
nondecreasing and $x\mapsto\min(0,x)$ is concave, $\beta(a,b)\ge\frac12(\beta(a,a)+\beta(b,b))$. So every entry, in
the same class or not, satisfies $v_p(G^e_{mm'})\ge\rho_m+\rho_{m'}$ with $\rho_m:=r_m+\frac12\beta(r_m,r_m)\le r_m$.
Lemma~\ref{lem:leibniz} gives $v_p(\det G^e)\ge2\sum_m\rho_m$, and therefore
\[
  v_p(\Delta^{\cR}_h)\ge\sum_m\beta(r_m,r_m)=-\sum_m(2r_m+1+M_m-n_m)_+=-\cb(h,p).\qedhere
\]
\end{proof}

\subsection{The class count}
Let $\varphi_1$ be the continuous, piecewise linear, nonincreasing function on $[2/5,\infty)$ with
\[
  \varphi_1(x)=\begin{cases}3-\tfrac52x,&\tfrac25\le x\le\tfrac12,\\ 2-\tfrac12x,&\tfrac12\le x\le\tfrac23,\\
  3-2x,&\tfrac23\le x\le1,\\ 2-x,&1\le x\le2,\end{cases}
\]
and $\varphi_1(x)=0$ for $x\ge2$. Then $\int_{2/5}^{2}\varphi_1=\frac3{16}+\frac{41}{144}+\frac49+\frac12=\frac{17}{12}$.

\begin{lemma}[Lemma~8A, class count]\label{lem:classcount}
Let $h\ge2$ and let $d$ be odd. If $2h\le5d$ and $d<2h$ then $\cb(h,d)\le h\,\varphi_1(d/h)$; more precisely
$\cb(h,d)=h\varphi_1(d/h)-\frac32$ if $d\le(2h-1)/3$ and $\cb(h,d)=h\varphi_1(d/h)$ otherwise. If $d\ge2h$ then $\cb(h,d)=0$.
(The formal version assumes $h\ge13$, which is needed only in Theorem~\ref{thm:generic}.)
\end{lemma}

\begin{proof}
Write $K=h$ and $\kappa=(d-1)/2$. If $d\ge2K$, every atom $l<K$ has $2l+1<d$, so all weights vanish, every class is a
singleton and every term is $(1+0-1)_+=0$. Now let $2K\le5d$ and $d<2K$; as $d$ is odd, $5d\ge2K+1$ and $3d\ne2K$. The
odd multiples of $d$ below $2K$ are $d$ and possibly $3d$. Hence an atom $l$ has weight $0$ if $l<\kappa$ (we call it
\emph{small}), weight $1$ if $l=\kappa$ or $l=\kappa+d$, and weight $2$ otherwise. Every atom after the first in its
class is $\ge d>\kappa$, so only the first atom $c$ of a class can be small. We evaluate $\cb$ class by class; a class
$\{c,c+d,\dots\}$ of size $n$ contributes $\sum_r(2r+1+M-n)_+$ over its atoms.
\begin{itemize}
\item[(IV)] $K<d<2K$. All classes are singletons, contributing their weights: $2(K-\kappa)-1=2K-d$.
\item[(III)] $(2K+1)/3\le d\le K$. Classes $c<K-d$ are pairs, with $c<\kappa$ small; they contribute $0+(1+2)=3$.
  The other classes $c\in[K-d,d-1]$ are singletons; those with $c\ge\kappa$ ($(d+1)/2$ of them, one of weight $1$)
  contribute $d$ in total. The atom $\kappa+d$ is $\ge K$. Total $3(K-d)+d=3K-2d$.
\item[(II)] $K/2\le d\le(2K-1)/3$. Classes $c<K-d$ are pairs and $c\in[K-d,d-1]$ singletons, all of weight $2$ (as
  $\kappa<K-d$). A pair contributes $(w(c)-1)_++1+M$: this is $3$ for $c<\kappa$, $2$ for $c=\kappa$ (weights
  $1,1$), and $4$ for $\kappa<c<K-d$. Total $3\kappa+2+4(K-d-\kappa-1)+2(2d-K)=2K-\frac d2-\frac32$.
\item[(I)] $(2K+1)/5\le d\le K/2$. Classes $c<K-2d$ are triples, with $c<\kappa$ (as $K-2d\le\kappa$), contributing
  $0+2+4=6$. Classes $c\in[K-2d,d-1]$ are pairs, contributing $3$ for $c<\kappa$ ($\frac{5d-1}2-K$ of them), $2$ for
  $c=\kappa$, and $4$ for $c>\kappa$ ($\frac{d-1}2$ of them). Total $3K-\frac52d-\frac32$.
\end{itemize}
In the four regions $K\varphi_1(d/K)$ equals $2K-d$, $3K-2d$, $2K-\frac d2$ and $3K-\frac52d$.
\end{proof}

\subsection{The generic primes}
For odd $p$ put
\[
  e(h,p):=\begin{cases}\cb(h,p),& p^2>2h\text{ and }5p\ge2h,\\ 2h\lev(p),&\text{otherwise.}\end{cases}
\]

\begin{theorem}[generic primes]\label{thm:generic}
For every odd prime $p$ and every $h$, $\|\Delta^{\cR}_h\|_p\le p^{e(h,p)}$. For every $\eps>0$ and all large $h$,
\[
  \sum_{p\text{ odd}}e(h,p)\log p\le\Bigl(\frac{133}{60}+\eps\Bigr)h^2 .
\]
\end{theorem}

\begin{proof}
The first statement is Proposition~\ref{prop:classbound} or Lemma~\ref{lem:nine}. For the sum let $h\ge13$. If
$5p\ge2h$ then $p^2\ge4h^2/25>2h$, so $e(h,p)=\cb(h,p)$, which vanishes for $p\ge2h$ and is at most $h\varphi_1(p/h)$
for $2h/5\le p<2h$ (Lemma~\ref{lem:classcount}). So the sum is at most
\[
  \sum_{p<2h/5}2h\lev(p)\log p+\sum_{2h/5\le p<2h}h\,\varphi_1(p/h)\log p .
\]
For the first sum, $\lev(p)\le1$ when $p^2>2h-1$ and $\lev(p)\log p\le\log(2h)$ always, so it is at most
$2h\bigl(\theta(2h/5)+(\sqrt{2h}+1)\log(2h)\bigr)=(\frac45+o(1))h^2$ by the prime number theorem $\theta(x)\sim x$.
For the second, cut $[2/5,2]$ into $M$ equal cells. On each cell $\varphi_1$ is at most its value at the left end,
and $\sum_{p\in hI}\log p=h|I|+o(h)$ for each fixed cell $I$. Hence the second sum is at most
$h^2\bigl(\int_{2/5}^2\varphi_1+\frac{16}{5M}\bigr)+o(h^2)$, and $M$ is arbitrary.
\end{proof}

% ===========================================================================
\section{\texorpdfstring{Denominators II: the primes $q$, $2$ and the divisors of $q^2-1$}{Denominators II: the primes q, 2 and the divisors of q^2-1}}\label{sec:arith2}
% ===========================================================================

\subsection{The pole order at \texorpdfstring{$q^2=1$}{q^2=1}}
For $f\in\cR$ let $\Ed f$ be its degree in $E$. Put
\begin{gather*}
  \rho_i:=(i+1-\lceil h/2\rceil)_+,\qquad \kappa_j:=(j-\lfloor h/2\rfloor)_+,\\
  \beta_h:=\sum_{i<h}\rho_i+\sum_{j<h}\kappa_j=\lfloor h^2/4\rfloor\le\frac{h^2+2h}4 .
\end{gather*}

\begin{lemma}[Lemma~5]\label{lem:poleorder}
Every coefficient of $\Delta^{\cR}_h\in\cR[X]$ has $E$-degree at most $\beta_h$.
\end{lemma}

\begin{proof}
In the monomial basis, the entry $(i,j)$ is $\mu^{\cR}(R)$ plus pole terms that do not involve $E$, where $R$ is the
quotient of $u^{i+j+1}$ by $D$, of degree $i+j+1-h$ (and $R=0$ if $i+j+1<h$). In $\mu^{\cR}(R)=\sum_{b<\deg R}\Delta^b[R'](0)\Omega_b$
each $\Omega_b$ has $E$-degree $b+1\le\deg R$. So the entry has $E$-degree at most $(i+j+1-h)_+\le\rho_i+\kappa_j$, and
the Leibniz expansion gives the claim. The value of $\beta_h$ is an elementary sum.
\end{proof}

\emph{Clearing $E$.} For $f=\sum c_{a,b}q^aE^b\in\cR$ with $\Ed f\le\beta$ put
$\operatorname{cl}_\beta(f):=\sum c_{a,b}\,q^a(q^2-1)^{\beta-b}\in\QQ[q]$. Since $(q^2-1)^k$ has integer coefficients,
$\|\operatorname{cl}_\beta f\|_p\le\|f\|_p$ for every $p$. If $v_1(v_0^2-1)=1$, then evaluating at $q=v_0$, $E=v_1$ gives
\begin{equation}\label{eq:clear}
  \operatorname{ev}\bigl(\operatorname{cl}_\beta f\bigr)=(v_0^2-1)^\beta\,\operatorname{ev}(f).
\end{equation}

\subsection{The \texorpdfstring{$q$}{q}-adic gain}
Let $\psi\colon\cR\to\QQ[[q]]$ be the ring homomorphism with $q\mapsto q$ and $E\mapsto1/(q^2-1)=-\sum_{n\ge0}q^{2n}$. Put
\[
  a_h:=\max\Bigl(0,\Bigl\lfloor\frac{h(h-1)-6h}{3}\Bigr\rfloor\Bigr)\ge\frac{h^2-7h-2}{3} .
\]

\begin{lemma}[Lemma~4$'$]\label{lem:qadic}
Every coefficient of $\psi(\Delta^{\cR}_h)\in\QQ[[q]][X]$ is divisible by $q^{a_h}$. Consequently, for every $k$, all
monomials of $\operatorname{cl}_{\beta_h}\bigl([X^k]\Delta^{\cR}_h\bigr)\in\QQ[q]$ have degree at least $a_h$.
\end{lemma}

\begin{proof}
\emph{Reflection.} From $\psi(\Omega_b)=(-1)^{b+1}q^2(1-q^2)^{-b-1}$ for $b\ge1$ and $\psi(\Omega_0)=-(1-q^2)^{-1}$, the
coefficient of $q^{2m}$ in $\psi(\Omega_b)$ is $-(-1)^b\binom{m+b-1}{b}$, and odd powers do not occur. Newton's backward
formula $f(-m)=\sum_b(-1)^b\binom{m+b-1}b\Delta^bf(0)$ for polynomials $f$ then gives
\begin{equation}\label{eq:reflect}
  \psi\bigl(\mu^{\cR}(R)\bigr)=-\sum_{m\ge0}q^{2m}R'(-m)\qquad(R\in\QQ[u]).
\end{equation}
\emph{Recombination.} For a polynomial $g$ write $g\varphi=R+\sum_kc_k/(u-x_k)$, $c_k=g(x_k)\omega_k$. Since
$\bigl((u-x_k)^{-1}\bigr)'(-m)=-4/(2k+1-2m)^2$ and
\[
  \sum_{m\ge0}\frac{q^{2m}}{(2k+1-2m)^2}=S_k(q)+q^{2k+1}\chi_2(q),
\]
where $\chi_2(q)=\sum_{j\ge0}q^{2j+1}/(2j+1)^2\in\QQ[[q]]$ (the terms $m\le k$ give $S_k$, the terms $m>k$ give $q^{2k+1}\chi_2(q)$), \eqref{eq:reflect} and Lemma~\ref{lem:pole}
combine to
\begin{equation}\label{eq:recomb}
  \psi\bigl(\Lambda^{\cR}_X(g)\bigr)=-\sum_{m\ge0}q^{2m}(g\varphi)'(-m)-4\,(X+\chi_2(q))\sum_{k<h}g(x_k)\,\omega_k\,q^{2k+1}.
\end{equation}
\emph{Slots.} For $g=u^{i+j}$ we have
$(g\varphi)'(-m)=\varphi'(-m)\,t^it^j+\varphi(-m)\bigl(it^{i-1}\bigr)t^j+\varphi(-m)\,t^i\bigl(jt^{j-1}\bigr)$ at $t=-m$, and
$g(x_k)=x_k^ix_k^j$. So the matrix $\psi(H)$ is a sum of rank-one matrices $c\,xy^{\mathsf T}$, where the row vector $x$ is
either a \emph{value row} $(t^i)_i$ or a \emph{derivative row} $(it^{i-1})_i$ at a node $t=-m$, or a pole row
$(x_k^i)_i$. A node $-m$ carries the factor $q^{2m}$ and a pole $k$ carries $q^{2k+1}$. Give these row types the
\emph{keys} $(m,\mathrm{val})$, $(m,\mathrm{der})$ and $k$, with \emph{prices} $2m$, $2m$ and $2k+1$.

\emph{Truncation and Andreief.} Let $M:=a_h+1$ and let $H_M$ be the same sum restricted to the nodes $m<M$. Then
$\psi(H)\equiv H_M$ entrywise modulo $q^{2M}$, so $\psi(\Delta^{\cR}_h)=\det\psi(H)\equiv\det H_M\pmod{q^{2M}}$. For a finite
sum of rank-one matrices, the Andreief (Cauchy--Binet) identity \cite{Andreief,Forrester2019} reads
\begin{equation}\label{eq:andreief}
  h!\,\det\Bigl(\sum_\iota c_\iota x_\iota y_\iota^{\mathsf T}\Bigr)
  =\sum_{f\colon[h]\to I}\Bigl(\prod_ic_{f(i)}\Bigr)\det\bigl(x_{f(i)}(j)\bigr)_{i,j}\det\bigl(y_{f(i)}(j)\bigr)_{i,j}.
\end{equation}
Here the rank-one terms are indexed by $\iota\in I$; the row vector $x_\iota$ and the power of $q$ in $c_\iota$ depend only
on the key of $\iota$ (the two terms $\varphi'(-m)\,t^it^j$ and $\varphi(-m)\,t^i(jt^{j-1})$ of a node share the key
$(m,\mathrm{val})$). A term of \eqref{eq:andreief} vanishes unless the rows $x_{f(i)}$ are pairwise distinct, so unless
the $h$ keys of $f(0),\dots,f(h-1)$ are pairwise distinct. The term is then divisible by $q$ to the sum of their prices. The map $(m,\mathrm{val})\mapsto3m$, $(m,\mathrm{der})\mapsto3m+2$,
$k\mapsto3k+1$ is injective and satisfies $2\cdot(\text{image})\le3\cdot\text{price}+6$. The images of $h$ distinct keys
are $h$ distinct naturals, with sum $\ge h(h-1)/2$; hence the prices sum to at least $(h(h-1)-6h)/3\ge a_h$. Since $h!$
is invertible and $2M>a_h$, $q^{a_h}$ divides every coefficient of $\psi(\Delta^{\cR}_h)$.

\emph{Clearing.} Let $f=[X^k]\Delta^{\cR}_h$. By Lemma~\ref{lem:poleorder} and \eqref{eq:clear} (with $v=(q,1/(q^2-1))$ in
$\QQ[[q]]$), $\operatorname{cl}_{\beta_h}f=(q^2-1)^{\beta_h}\psi(f)$ in $\QQ[[q]]$, which is divisible by $q^{a_h}$. As
$\operatorname{cl}_{\beta_h}f$ is a polynomial in $q$, its monomials of degree $<a_h$ vanish.
\end{proof}

\subsection{The prime 2}

\begin{lemma}[Lemma~10]\label{lem:two}
Let $q\in\{3,5\}$. For every $h\ge1$ and $k\ge0$, $\bigl|[X^k]\Delta_h\bigr|_2\le2^{\beta_h+h(\lfloor\log_2h\rfloor+6)}$.
\end{lemma}

The only property of $q$ used is $v_2(q^2-1)=3$. We work directly over $\QQ$ with the anchored basis, for which
$\det(\Lambda_X(a_ib_j))=\Delta_h(X)$. Put $L:=\lfloor\log_2h\rfloor$ and write $|\cdot|_2$ coefficientwise on $\QQ[X]$.

\begin{lemma}\label{lem:fdseq}
Let $T\colon\ZZ_{\ge0}\to\QQ$ satisfy $T(l+1)=Q\,T(l)+e\,(s+2l)^{-2}$ with $v_2(Q-1)\ge2$, $v_2(e)\ge0$, $v_2(T(0))\ge0$
and $s$ odd. Then $v_2(\Delta^mT(0))\ge v_2(m!)+m-2-\lfloor\log_2m\rfloor$ for $m\ge1$.
\end{lemma}

\begin{proof}
Put $\eta=Q-1$ and $g(l)=(s+2l)^{-2}$. Then $\Delta T=\eta T+eg$, and iterating,
$\Delta^mT(0)=\eta^mT(0)+e\sum_{k<m}\eta^{m-1-k}\Delta^kg(0)$. By induction on $k$,
$\Delta^kg(0)=(-2)^kk!\,H_k(s)/P_k(s)$ with $P_k(s)=\prod_{a\le k}(s+2a)$ odd and $H_k(s)=\sum_{a\le k}(s+2a)^{-1}$
$2$-integral, so $v_2(\Delta^kg(0))\ge k+v_2(k!)$. Using $k-1-\lfloor\log_2k\rfloor\le v_2(k!)$ and $v_2(m!)\le m-1$, each
term has valuation at least $2(m-1-k)+2k-1-\lfloor\log_2k\rfloor\ge v_2(m!)+m-2-\lfloor\log_2m\rfloor$, and
$v_2(\eta^mT(0))\ge2m$ is larger.
\end{proof}

\begin{proof}[Proof of Lemma~\ref{lem:two}]
We show that each entry satisfies
\begin{equation}\label{eq:twoentry}
  v_2\bigl(\Lambda_X(a_ib_j)\bigr)\ge h-(i+j+1)-(i+j+1-h)_+-L-6 .
\end{equation}
\emph{Entries with $\pi:=i+j+1-h\ge1$.} Then $a_ib_j=D\cdot P$ with $P=\prod_{h-j\le m<i}(u-x_m)$, so $a_ib_j\varphi=uP$
is a polynomial: $B(a_ib_j)=0$ and $A(a_ib_j)=\mu(uP)$. We have $uP=2^{-\pi}H(2u)$ with $H\in\ZZ[T]$ of degree $\pi$.
Since $\Delta^b[(2t)^n](0)=2^nb!\,S(n,b)$ (Stirling numbers of the second kind) vanishes for $n<b$,
$v_2(\Delta^b[(uP)'](0))\ge1-\pi+b+v_2(b!)$. And $v_2(\Omega_b)=-3(b+1)$. By \eqref{eq:newton} and
$v_2(b!)\ge b-1-\lfloor\log_2b\rfloor$,
\[
  v_2(A(a_ib_j))\ge\min_{b<\pi}\bigl(-\pi-2-2b+v_2(b!)\bigr)\ge-2\pi-2-\lfloor\log_2\pi\rfloor ,
\]
which implies \eqref{eq:twoentry}.

\emph{Entries with $N:=h-1-i-j\ge0$.} Then $a_ib_j\varphi=u/\prod_{l=0}^{N}(u-\xi_l)$ with $\xi_l:=x_{i+l}=-c/2-l$,
$c:=2i+1$. Since $\xi_l-\xi_m=m-l$, its partial fractions are a constant (which $\mu$ kills) plus
$\sum_l(-1)^l\binom Nl\xi_l/(N!\,(u-\xi_l))$. By the proof of Lemma~\ref{lem:pole}, both $A(a_ib_j)$ and $B(a_ib_j)$ have the form
$\pm\frac4{N!}\sum_l(-1)^l\binom Nl\xi_lT(l)=\pm\frac4{N!}\Delta^N[\xi_lT(l)](0)$, with $T(l)=S_{i+l}(q)$ or
$T(l)=q^{2(i+l)+1}$. Both satisfy the recursion of Lemma~\ref{lem:fdseq} with $Q=q^2$: $S_{i+l+1}=q^2S_{i+l}+(2i+3+2l)^{-2}$
($e=1$, $s=2i+3$), and $q^{2(i+l+1)+1}=q^2\cdot q^{2(i+l)+1}$ ($e=0$). From $\Delta^N[l\,T(l)](0)=N\,\Delta^{N-1}[T(\cdot+1)](0)$, and
since $T(\cdot+1)$ satisfies a recursion of the same kind,
\begin{multline*}
  v_2\Bigl(\frac4{N!}\Delta^N[\xi_lT(l)](0)\Bigr)\ge2-v_2(N!)\\
  +\min\Bigl(-1+v_2(N!)+N-2-\lfloor\log_2N\rfloor,\ v_2(N!)+N-3-\lfloor\log_2(N-1)\rfloor\Bigr),
\end{multline*}
for $N\ge2$, which is at least $N-1-\lfloor\log_2N\rfloor$. For $N\le1$ the entry is $\pm4\,\xi_0T(0)$ or
$\pm4\bigl(\xi_0\Delta T(0)-T(1)\bigr)$, of valuation $\ge1$, since $v_2(\xi_0)=-1$ and $T(0),T(1),\Delta T(0)$ are
$2$-integral. This implies \eqref{eq:twoentry}.

\emph{Assembly.} Since $(i+j+1-h)_+\le\rho_i+\kappa_j$, \eqref{eq:twoentry} gives
$|\Lambda_X(a_ib_j)|_2\le\bigl(2^{i+1+\rho_i+L+6-h}\bigr)\bigl(2^{j+\kappa_j}\bigr)$. The Leibniz expansion, now over
$\QQ[X]$ with the coefficientwise $2$-adic norm, gives
$|[X^k]\Delta_h|_2\le2^{\sum_i(i+1)+\sum_jj+\beta_h+h(L+6)-h^2}=2^{\beta_h+h(L+6)}$.
\end{proof}

\subsection{The denominator}

\begin{theorem}[arithmetic bound]\label{thm:arith}
Let $q\in\{3,5\}$. There are rationals $d_h>0$ with $d_h\Delta_h\in\ZZ[X]$ for every $h\ge1$, such that for every
$\eps>0$ and all large $h$,
\[
  \log d_h\le(\alpha_q+\eps)h^2,\qquad
  \alpha_q=\frac{133}{60}+\frac{\log2}4+\ind{q=5}\frac{\log3}4-\frac{\log q}3 .
\]
\end{theorem}

\begin{proof}
Let $c$ be a coefficient of $\Delta_h$; then $c=\operatorname{ev}_q(f)$ with $f$ the corresponding coefficient of
$\Delta^{\cR}_h$. We bound $|c|_p$ for every prime $p$.
\begin{itemize}
\item $p=2$: $|c|_2\le2^{\beta_h+h(L+6)}$ by Lemma~\ref{lem:two}.
\item $p=q$: by \eqref{eq:clear}, $c=\operatorname{ev}(\operatorname{cl}_{\beta_h}f)/(q^2-1)^{\beta_h}$, and $q^2-1$ is a $q$-unit.
  By Lemma~\ref{lem:qadic} every monomial of $\operatorname{cl}_{\beta_h}f$ has degree $\ge a_h$, so
  $|c|_q\le\|\operatorname{cl}_{\beta_h}f\|_q\,q^{-a_h}\le\|f\|_q\,q^{-a_h}\le q^{2h\lev(q)-a_h}$ by Lemma~\ref{lem:nine}.
\item $p$ odd with $p\mid q^2-1$ (so $p=3$, $q=5$): by Lemma~\ref{lem:poleorder} and \eqref{eq:clear},
  $|c|_p=|\operatorname{ev}(\operatorname{cl}_{\beta_h}f)|_p\,|q^2-1|_p^{-\beta_h}\le\|f\|_p\,|q^2-1|_p^{-\beta_h}\le
  p^{2h\lev(p)+v_p(q^2-1)\beta_h}$, by Lemma~\ref{lem:nine}.
\item every other $p$: $q$ and $1/(q^2-1)$ are $p$-integral, so $|c|_p\le\|f\|_p\le p^{e(h,p)}$ by
  Theorem~\ref{thm:generic}.
\end{itemize}
Let $N_h:=\prod_pp^{e_p}$, where $e_2=\beta_h+h(L+6)$, $e_q=2h\lev(q)$, $e_3=2h\lev(3)+\beta_h$ if $q=5$, and $e_p=e(h,p)$
otherwise. Only primes $p<2h+q^2$ occur, because $e(h,p)=0$ for $p\ge2h$. Put $d_h:=N_h/q^{a_h}$. Then
$|d_hc|_p\le1$ for every prime $p$, so $d_hc\in\ZZ$. Finally
\begin{multline*}
  \log d_h\le\sum_{p\nmid2q(q^2-1)}e(h,p)\log p+\Bigl(\frac{h^2+2h}4+h(L+6)\Bigr)\log2\\
  +\ind{q=5}\frac{h^2+2h}4\log3-a_h\log q+O(h\log h),
\end{multline*}
where the $O(h\log h)$ collects the terms $2h\lev(p)\log p\le2h\log(2h)$ for $p\in\{3,q\}$. Theorem~\ref{thm:generic}
and $a_h\ge(h^2-7h-2)/3$ give the claim.
\end{proof}

% ===========================================================================
\section{The size of \texorpdfstring{$\Delta_h(\chi)$}{Delta_h(chi)}: Theorem B}\label{sec:analytic}
% ===========================================================================

\subsection{Dual certificates}
Let $\psi(u):=u\log|u|-u$ (with $\psi(0)=0$), an antiderivative of $\log|u|$, so that
$\int_\alpha^\beta\log|x-y|\,dy=\psi(x-\alpha)-\psi(x-\beta)$. Put
\[
  V(x):=x\log w-\psi(x+1)+\psi(x)=x\log w-(x+1)\log(x+1)+x\log x+1
\]
for $x\ge0$.
A \emph{step density} is $\mu=\sum_{j<J}c_j\mathbf 1_{[a_j,a_{j+1})}$ with $0=a_0\le a_1\le\dots\le a_J$, $0\le c_j\le2$ and
$\int\mu=1$. Its potential and energy are
\[
  U^\mu(x):=\int\log|x-y|\,\mu(y)\,dy,\qquad E(\mu):=\iint\log|x-y|\,\mu(x)\mu(y)\,dx\,dy .
\]

\begin{definition}\label{def:cert}
A \emph{dual certificate} for $q$ consists of a step density $\mu$, a real number $\lambda$, integers $L,X,N\ge1$ with
$a_J\le N\le X$, and real numbers $S_k$ ($0\le k<XL$) such that
\begin{enumerate}
\item $F(x):=2U^\mu(x)+V(x)\le S_k$ for $x\in[k/L,(k+1)/L]$, for every $k<XL$;
\item $1/X+\log w<0$ and $\log X+X\log w<\lambda$.
\end{enumerate}
Its \emph{exponent} is $\Phi:=-E(\mu)+\lambda+\frac2L\sum_{k<XL}(S_k-\lambda)_+$.
\end{definition}

\begin{theorem}[Theorem~B]\label{thm:B}
For every dual certificate for $q$, as $h\to\infty$,
\[
  \log|\Delta_h(\chi_2(1/q))|\le\Phi h^2+O(h\log h).
\]
\end{theorem}

The constant in $O(h\log h)$ depends on the certificate. The heuristic behind the theorem: by \eqref{eq:value},
$\Delta_h(\chi)$ is a Hankel determinant of a positive-weight functional on the nodes $t=1,2,\dots$ with the
derivative structure of $(g\varphi)'$, i.e.\ a discrete log-gas in the external field $V$ with at most two
particles per node. The dual certificate is a feasible point of the dual of the corresponding constrained
equilibrium problem \cite[Ch.~I]{SaffTotik}, and the proof below is a discrete version of weak duality.

\subsection{Node expansion}
We have $\Lambda(u^n)=\sum_{s\ge1}w^s\bigl(ns^{n-1}\varphi(s)+s^n\varphi'(s)\bigr)$. Let $H^{(T)}$ be the matrix of
truncated sums over $1\le s\le T$. Then $H^{(T)}\to(\Lambda(u^{i+j}))_{i,j}$ entrywise, so by continuity of the
determinant $|\Delta_h(\chi)|\le\sup_T|\det H^{(T)}|$. Expanding $(i+j)s^{i+j-1}=(is^{i-1})s^j+s^i(js^{j-1})$,
\[
  H^{(T)}=\sum_{s\le T}\Bigl(w^s\varphi'(s)\,v_sv_s^{\mathsf T}+w^s\varphi(s)\,v'_sv_s^{\mathsf T}+w^s\varphi(s)\,v_sv_s'^{\mathsf T}\Bigr),
\]
with the value row $v_s=(s^j)_{j<h}$ and the derivative row $v'_s=(js^{j-1})_{j<h}$. We call $(s,\mathrm{val})$ and
$(s,\mathrm{der})$ the two \emph{slots} of the node $s$. Since $\varphi'/\varphi=1/s-\sum_k1/(s-x_k)$ is a difference of two
numbers in $[0,h]$ for $s\ge1$, all three coefficients are bounded by
\[
  b(s):=w^s\,h\,\varphi(s).
\]

\begin{lemma}[minors with derivative rows]\label{lem:slotdet}
Let $h\ge2$ and let $(s_i,\tau_i)_{i<h}$ be slots at integer nodes $s_i\ge0$. Then
\[
  \bigl|\det(\text{row of slot }(s_i,\tau_i))_{i<h}\bigr|\le h^h\Bigl(1+\frac2h\Bigr)^{h^2}\prod_{i<j}\operatorname{pf}(s_i,s_j),
\]
where $\operatorname{pf}(a,b)=|a-b|$ for $a\ne b$ and $\operatorname{pf}(a,a)=1$.
\end{lemma}

\begin{proof}
Let $\zeta$ be a primitive $h$-th root of unity and $r=1/h$. For $j<h$,
$\sum_{l<h}\zeta^{l(h-1)}(t+r\zeta^l)^j=h\,r\,jt^{j-1}$ (only the binomial term with $h\mid h-1+m$, i.e.\ $m=1$, survives;
this needs $h\ge2$). Since $hr=1$, the derivative row at $t$ is the sum
$\sum_l\zeta^{l(h-1)}(\text{value row at }t+r\zeta^l)$. By multilinearity the determinant is a sum of at most
$h^{\#\mathrm{der}}$ Vandermonde determinants at points $z_i$ with $|z_i-s_i|\le r$, each with a coefficient of modulus
$1$. Since the nodes are integers, each Vandermonde factor satisfies
$|z_j-z_i|\le|s_j-s_i|+2r\le(1+2r)\operatorname{pf}(s_i,s_j)$.
So the determinant is at most $h^h(1+2/h)^{h(h-1)/2}\prod_{i<j}\operatorname{pf}(s_i,s_j)$.
\end{proof}

Apply \eqref{eq:andreief} to $H^{(T)}$, indexing the rank-one terms by triples (node, type). A term with two equal
row vectors vanishes, so only assignments $f$ whose \emph{x-slots} are distinct contribute; in particular each node
carries at most two of the $h$ indices. With Lemma~\ref{lem:slotdet} for both determinants,
\begin{equation}\label{eq:pattern}
  |\det H^{(T)}|\le\frac{K_h^2}{h!}\sum_{f}\ \prod_{i<h}b(s_i)\prod_{i<j}\operatorname{pf}(s_i,s_j)^2,\qquad K_h:=h^h(1+2/h)^{h^2},
\end{equation}
where $f$ runs over the assignments $i\mapsto(s_i,\tau_i)$ with distinct x-slots and $s_i$ is the node of $f(i)$.

\subsection{A lattice energy inequality}

\begin{lemma}[P\'olya-type inequality]\label{lem:polya}
For $M\ge1$ and real $\nu_0,\dots,\nu_{M-1}$,
\[
  \sum_{i\ne j}\nu_i\nu_j\log|i-j|\le\log2\sum_i\nu_i^2+\log M\Bigl(\sum_i\nu_i\Bigr)^2 .
\]
\end{lemma}

\begin{proof}
For $M=1$ the claim reads $0\le\log2\,\nu_0^2$. Let $M\ge2$. We follow P\'olya's criterion for positive definite
functions \cite{Polya1949} in a finite form. Let $f(0)=\log M+\log2$, $f(d)=\log(M/d)$ for $1\le d\le M$, and $f(d)=0$
for $d>M$. Then $f$ is convex on
$\ZZ_{\ge0}$ (this uses $\log(k-1)+\log(k+1)\le2\log k$ for $2\le k<M$, $f(M-1)-2f(M)+f(M+1)=\log\frac M{M-1}\ge0$,
and $f(0)-2f(1)+f(2)=0$ when $M\ge2$).
Hence $f(d)=\sum_{k=1}^{M}\Delta^2f(k-1)\,(k-d)_+$ with $\Delta^2f\ge0$. Each triangle kernel $(k-|i-j|)_+$ is the
number of windows of length $k$ containing both $i$ and $j$, hence positive semidefinite. So
$\sum_{i,j}\nu_i\nu_jf(|i-j|)\ge0$, which rearranges to the claim.
\end{proof}

For a reference sequence $m=(m_s)_{s\ge0}$ with finite support, $\sum_sm_s=h$, put
\[
  P_m(s):=\sum_{s'\ne s}m_{s'}\log|s-s'|,\qquad E_m:=\sum_{s\ne s'}m_sm_{s'}\log|s-s'|.
\]

\begin{lemma}[discrete weak duality]\label{lem:dual}
Let $m$ be as above with $\sum_sm_s^2\le2h$, and let $\lambda'\in\RR$. Put $g(s):=\log b(s)+2P_m(s)-\lambda'$. For every assignment $f$ in \eqref{eq:pattern},
\[
  \prod_ib(s_i)\prod_{i<j}\operatorname{pf}(s_i,s_j)^2\le\exp\Bigl(-E_m+\lambda'h+8h\log2+\sum_ig(s_i)\Bigr).
\]
\end{lemma}

\begin{proof}
Let $n_s:=\#\{i:s_i=s\}\le2$, so $\sum_sn_s=h$ and $\sum_sn_s^2\le2h$. Pairs at the same node contribute
$\operatorname{pf}=1$, so $2\sum_{i<j}\log\operatorname{pf}(s_i,s_j)=\sum_{s\ne s'}n_sn_{s'}\log|s-s'|$. With $\nu:=n-m$ (total mass $0$),
\[
  \sum_{s\ne s'}n_sn_{s'}\log|s-s'|=-E_m+2\sum_sn_sP_m(s)+\sum_{s\ne s'}\nu_s\nu_{s'}\log|s-s'| ,
\]
and by Lemma~\ref{lem:polya} the last sum is at most $\log2\sum\nu_s^2\le2\log2(\sum n_s^2+\sum m_s^2)\le8h\log2$.
Finally $\sum_sn_sP_m(s)=\sum_iP_m(s_i)$.
\end{proof}

\subsection{Discretization}
Given a step density $\mu$ supported in $[0,N]$, let $m_s:=h\int_{(s-1)/h}^{s/h}\mu$ for $s\ge1$ and $m_0:=0$ (cell
averages). Then $\sum m_s=h$, $0\le m_s\le2$ (so $\sum m_s^2\le2h$), and $m_s=0$ for $s>Nh$.

\begin{lemma}\label{lem:discrete}
With these $m_s$, for all $h\ge2$ and $s\ge1$:
\begin{enumerate}
\item $\log b(s)\le\log h+\log s-h\log h+h\,V(s/h)$;
\item $P_m(s)\le h\,U^\mu(s/h)+h\log h+4+2\log(Nh+1)$;
\item $E_m\ge h^2E(\mu)+h^2\log h-4h\bigl(1+\log(Nh+1)\bigr)$.
\end{enumerate}
\end{lemma}

\begin{proof}
(i) By concavity of $\log$, $\log(c+\frac12)\ge\int_c^{c+1}\log y\,dy=\psi(c+1)-\psi(c)$ for $c\ge0$, so
$\log D(s)=\sum_{k<h}\log(s+k+\frac12)\ge\psi(s+h)-\psi(s)=h\bigl(\psi(x+1)-\psi(x)\bigr)+h\log h$ with $x=s/h$. Insert this into
$\log b(s)=s\log w+\log h+\log s-\log D(s)$.

(ii) We have $h\,U^\mu(s/h)=\sum_{s'}\int_{I_{s'}}\log|s-u|\,\rho(u)\,du-h\log h$, where $I_{s'}=[s'-1,s']$ and
$\rho(u):=\mu(u/h)$; this density has mass $m_{s'}$ on $I_{s'}$, total mass $h$, and is bounded by $2$. For
$s'<s$ and $u\in I_{s'}$, $\log|s-u|\ge\log(s-s')$, so these cells contribute with the right sign. For $s'=s$ the
discrete term is $0$ and $\int_{I_s}(-\log|s-u|)\,\rho\le2$. For $s'=s+d$, $d\ge1$, the deficit is at most
$2\int_{d-1}^{d}(\log d-\log u)\,du=2\bigl(1-(d-1)\log\tfrac d{d-1}\bigr)\le2/d$. Summing over $d\le Nh$ gives (ii).

(iii) Similarly $h^2E(\mu)+h^2\log h=\sum_{s,s'}\iint_{I_s\times I_{s'}}\log|u-u'|\,\rho(u)\rho(u')$. The diagonal cells
contribute $\le0$ because $|u-u'|\le1$ there. For $|s-s'|=d\ge1$ we have $\log|u-u'|\le\log(d+1)\le\log d+1/d$. Since $m_{s'}\le2$, summing
$m_sm_{s'}/d$ over the pairs gives at most $\sum_sm_s\cdot2\cdot2\sum_{d\le Nh}1/d\le4h(1+\log(Nh+1))$.
\end{proof}

\subsection{Proof of Theorem~\ref{thm:B}}
Fix a certificate and $h\ge2$. Take $m$ as above and $\lambda':=h\lambda+h\log h$. By Lemma~\ref{lem:discrete}, with
$\epsilon_h:=4+2\log(Nh+1)$ and $x=s/h$,
\begin{gather}
  g(s)\le h\bigl(F(x)-\lambda\bigr)+\log h+\log s+2\epsilon_h,\label{eq:gbound}\\
  -E_m+\lambda'h\le h^2\bigl(\lambda-E(\mu)\bigr)+O(h\log h).\notag
\end{gather}
\emph{Near and far nodes.} Let $S_0:=Xh$. Since each node carries at most two indices,
$\sum_ig(s_i)\le2\sum_{1\le s\le S_0}g(s)_++\sum_{i:\,s_i>S_0}g(s_i)$. Summing over $f$ in \eqref{eq:pattern} with
Lemma~\ref{lem:dual} therefore gives
\[
  |\det H^{(T)}|\le\frac{K_h^2}{h!}\,e^{-E_m+\lambda'h+8h\log2+2\sum_{s\le S_0}g(s)_+}\Bigl(3S_0+3\sum_{s>S_0}e^{g(s)}\Bigr)^h .
\]
\emph{Near nodes.} Each $x=s/h\le X$ lies in a cell $[k/L,(k+1)/L]$ with $k<XL$, where $F\le S_k$; each cell contains at
most $h/L+1$ points $s/h$. By \eqref{eq:gbound},
$\sum_{s\le S_0}g(s)_+\le h(h/L+1)\sum_k(S_k-\lambda)_++S_0(\log h+\log S_0+2\epsilon_h)$.

\emph{Far nodes.} Put $\delta:=-\max\bigl(1/X+\log w,\ (\log X+X\log w-\lambda)/X\bigr)>0$. For $x\ge N$,
$U^\mu(x)\le\log x$, and $V(x)\le x\log w-\log x$ (because $1-x\log(1+1/x)\le1/(x+1)\le\log(1+1/x)$). Hence, using
$\log x\le\log X+(x-X)/X$, for $x\ge X$
\begin{align*}
  F(x)-\lambda&\le\log x+x\log w-\lambda\\
  &\le(\log X+X\log w-\lambda)+(x-X)(1/X+\log w)\le-\delta x .
\end{align*}
So $e^{g(s)}\le h\,s\,e^{2\epsilon_h}e^{-\delta s}\le\frac2\delta h\,e^{2\epsilon_h}e^{-\delta s/2}$ for $s>S_0$, and the
far sum is $e^{O(\log h)}$.

\emph{Assembly.} $\log(K_h^2/h!)=O(h\log h)$, and $\log\bigl(3S_0+3\sum_{s>S_0}e^{g(s)}\bigr)=O(\log h)$. Collecting,
\[
  \log|\det H^{(T)}|\le h^2\Bigl(-E(\mu)+\lambda+\frac2L\sum_k(S_k-\lambda)_+\Bigr)+O(h\log h)
\]
uniformly in $T$, and the theorem follows from $|\Delta_h(\chi)|\le\sup_T|\det H^{(T)}|$. \qed

\begin{remark}
Let $C(w):=\sup\{E(\rho)+\int V\rho\}$ over densities $\rho$ on $[0,\infty)$ with compact support, $0\le\rho\le2$ and
$\int\rho=1$ (the constrained equilibrium constant). For such $\rho$, $E(\rho-\mu)\le0$ \cite[Lemma~I.1.8]{SaffTotik},
so $E(\rho)+\int V\rho\le-E(\mu)+\int F\rho\le-E(\mu)+\lambda+2\int_0^X(F-\lambda)_+\le\Phi$, using $\rho\le2$, $F<\lambda$
beyond $X$ and $F\le S_k$ on the cells. Hence
$\Phi\ge C(w)$. Numerically $C(1/9)\approx-2.1003$ and
$C(1/25)\approx-2.4565$. The finite-$h$ values $h^{-2}\log|\Delta_h(\chi)|$, computed in ball arithmetic up to $h=320$,
approach these constants [numerical], so Theorem~\ref{thm:B} is essentially sharp for this construction. Only the
form with an explicit certificate is used.
\end{remark}

% ===========================================================================
\section{The two certificates}\label{sec:cert}
% ===========================================================================

\begin{theorem}\label{thm:cert}
There are dual certificates with exponent $\Phi<-2.078$ for $q=3$ and $\Phi<-2.385$ for $q=5$.
\end{theorem}

The certificates are explicit rational data (Table~\ref{tab:cert}). The density $\mu$ is constant on the $J$ cells
$[js/L,(j+1)s/L)$ with dyadic heights $c_j\in[0,2]$ (denominator $2^{16}$) and exact mass
$\sum_jc_j\,s/L=1$. The bounds $S_k$ are rationals. All checks are inequalities between explicit rational numbers
and finitely many values of $\log$ at rational points, and are verified as follows.
\begin{itemize}
\item \emph{Logarithms.} A table of dyadic intervals of width $2\cdot2^{-40}$ for $\log n$, $n\le n_{\max}$, each checked
  against the series of $\log(1-x)$ with an explicit tail bound, after reducing the argument by powers of $2$; $\log2$
  is enclosed to width $2^{-100}$.
\item \emph{Potential.} By summation by parts,
  $2U^\mu(x)+V(x)=\sum_jD_j\,\psi(x-\alpha_j/L)+x\log w$, with $D_j=2(c_j-c_{j-1})$ at $\alpha_j=js$, plus the terms
  $\psi(x)$ and $-\psi(x+1)$ of $V$. On each grid cell $[k/L,(k+1)/L]$ every term is bounded above by an affine
  function: by the chord where the term is convex, and by a tangent (at the midpoint, or at the endpoint away from
  $\alpha_j/L$) where it is concave. Summing the affine bounds and taking the larger endpoint value bounds $F$ on the
  cell. All values are rational enclosures.
\item \emph{Energy.} $E(\mu)=\sum_{j,j'}c_jc_{j'}\iint\log|x-y|$ over pairs of cells, and the double integral over
  $[a_1,b_1]\times[a_2,b_2]$ equals $-\Psi(b_1-b_2)+\Psi(b_1-a_2)+\Psi(a_1-b_2)-\Psi(a_1-a_2)$ with
  $\Psi(u)=\frac12u^2\log|u|-\frac34u^2$. A rational lower bound $E_{\mathrm{lo}}\le E(\mu)$ is computed in interval
  arithmetic.
\item \emph{Tail.} $1/X+\log w<0$ and $\log X+X\log w<\lambda$, from the enclosures of $\log X$ and $\log q$.
\end{itemize}
The certified exponent is the rational number $\Phi_{\mathbb Q}:=-E_{\mathrm{lo}}+\lambda+\frac2L\sum_k(S_k-\lambda)_+\ge\Phi$.

\begin{table}[ht]
\centering
\begin{tabular}{@{}lcc@{}}
\toprule
 & $q=3$ & $q=5$\\
\midrule
grid width $1/L$ & $1/100$ & $1/40$\\
near region $[0,X]$ (cells $XL$) & $[0,3]$ ($300$) & $[0,2]$ ($80$)\\
density: $J$ cells of width $s/L$ & $30\times\frac1{20}$ & $23\times\frac1{20}$\\
$N$; entries of the log table & $2$; $804$ & $2$; $244$\\
$-E_{\mathrm{lo}}$ & $1.532174204$ & $1.768055783$\\
$\lambda$ & $-3.900735042$ & $-4.649276378$\\
$\frac2L\sum_k(S_k-\lambda)_+$ & $0.290234795$ & $0.496069677$\\
$\Phi_{\mathbb Q}$ & $-2.0783260426\ldots$ & $-2.3851509182\ldots$\\
\bottomrule
\end{tabular}
\caption{The two dual certificates (decimal values rounded).}\label{tab:cert}
\end{table}

Exactly, $\Phi_{\mathbb Q}=-25979075532629/12500000000000$ for $q=3$ and $\Phi_{\mathbb Q}=-9540603672893/4000000000000$
for $q=5$.

The data are generated by a script that optimises the density in floating point, rounds it to dyadic rationals,
and then replicates the interval arithmetic exactly in rational arithmetic. The script is not trusted: in the
formalisation every check above is a decidable statement about the data, proved by evaluation in the Lean kernel
(Section~\ref{sec:lean}).

% ===========================================================================
\section{Nonvanishing: Lemma P}\label{sec:nonvan}
% ===========================================================================

\subsection{Statement}
For $r\in\{0,1,2\}$ put $c_r:=1/\prod_{a\in\{0,1,2\},a\ne r}(a-r)$, so $(c_0,c_1,c_2)=(\frac12,-1,\frac12)$. Recall
$T(M)=\sum_{s\le M,\,s\text{ odd}}z^s/s^2$. Define three $3\times3$ rational matrices, indexed by $r,s\in\{0,1,2\}$:
\begin{align*}
  M_S(q)_{rs}&=\frac{c_r}{1-w}+\ind{r=s}\,4q^{2r+1}T(2r-1),\\
  M_L(q)_{rs}&=\frac{w\,c_r}{1-w}+\ind{r=s}\,4q^{2r+1}T(2r+1),\\
  M_D(q)_{rs}&=w\,c_r\Bigl(\frac{e_{rs}}{1-w}+\frac{2w}{(1-w)^2}\Bigr)-\ind{r=s}\,4\bigl(r+\tfrac12\bigr)q^{2r+1}T(2r+1),
\end{align*}
where $e_{rs}=\frac{11}2-2r$ if $r=s$ and $e_{rs}=(3-r-s)+\frac52$ if $r\ne s$.

\begin{theorem}[Lemma~P]\label{thm:lemmaP}
Let $q\ge2$ and let $\ell\ge7$ be a prime with $\ell\nmid q(q^2-1)$ such that $\det M_S(q)$, $\det M_L(q)$ and
$\det M_D(q)$ are $\ell$-adic units. Then $\Delta_{3\ell}(a/b)\ne0$ for all $a\in\ZZ$ and $b\ge1$ with $\ell\nmid b$.
\end{theorem}

\begin{theorem}\label{thm:nonvan}
For $q\in\{3,5\}$ and every $a\in\ZZ$, $b\ge1$, we have $\Delta_h(a/b)\ne0$ for infinitely many $h$.
\end{theorem}

\begin{proof}
Exact evaluation gives
\[
\begin{array}{lccc}
 & \det M_S & \det M_L & \det M_D\\[2pt]
q=3: & 6642 & \frac{10896962}{225} & -\frac{5947389}{80}\\[4pt]
q=5: & \frac{3531250}{27} & \frac{246357730}{243} & -\frac{6934894375}{3888}
\end{array}
\]
(the factor $9q^2+1$ of $\det M_S=8q^6(9q^2+1)/(9(q^2-1))$ comes from $T(3)=z+z^3/9$). Every prime $\ell$ larger than
all these numerators and denominators, than $q^2$, and than $b$, satisfies the hypotheses of Theorem~\ref{thm:lemmaP}
and $\ell\nmid b$. There are infinitely many such primes, and each gives $h=3\ell$.
\end{proof}

The rest of this section proves Theorem~\ref{thm:lemmaP}. Fix $q,\ell$ as there and $h=3\ell$; then $\ell^2\ge7\ell>2h$.
Let $\ZZ_{(\ell)}$ be the ring of $\ell$-integral rationals and $\bar{\cdot}\colon\ZZ_{(\ell)}\to\FF_\ell$ the reduction. Put
$\kappa:=(\ell-1)/2$. The indices $i<h$ fall into the classes $C_\gamma=\{\gamma,\gamma+\ell,\gamma+2\ell\}$, $\gamma<\ell$;
we write $i=\gamma+r\ell$ with $r\in\{0,1,2\}$. The class $\gamma$ has \emph{type} S if $\gamma<\kappa$, D if $\gamma=\kappa$, and L if
$\gamma>\kappa$. Note that $\ell\mid2i+1$ iff $i\in C_\kappa$; put $\delta_i:=\ind{i\in C_\kappa}$.

\subsection{Steps 1--2: the Lagrange Gram matrix}
For $r_0\in\QQ$ let $\Lambda_{r_0}:=A+r_0B$ and $G(r_0)_{ij}:=\Lambda_{r_0}(f_if_j)$, with the Lagrange basis $f_i=D/(u-x_i)$. By
\eqref{eq:gram}, $\det G(r_0)=(\det C_f)^2\Delta_h(r_0)$. Put
\[
  \mathcal E(r_0)_{ij}:=\nu_i\,G(r_0)_{ij}\,\ell^{-\delta_j},\qquad \nu_i:=\ell^2/D'(x_i).
\]
Then $\det\mathcal E(r_0)\ne0$ implies $\Delta_h(r_0)\ne0$; no value of $\det C_f$ is needed. As in the proof of
Lemma~\ref{lem:nine}, for $i\ne j$ we have $G(r_0)_{ij}=\mu(F_{ij})$ with $F_{ij}=u\prod_{m\ne i,j}(u-x_m)$, and
$G(r_0)_{ii}=\mu(Q_i)+4x_iD'(x_i)\bigl(S_i(q)-r_0q^{2i+1}\bigr)$. Hence
\[
  \mathcal E(r_0)=\mathcal E(0)+r_0\operatorname{diag}(d_i),\qquad d_i=-4\ell^{2-\delta_i}x_iq^{2i+1},\qquad v_\ell(d_i)=2 .
\]
If $\mathcal E(0)$ is $\ell$-integral and $\ell\nmid b$, then $\mathcal E(a/b)$ is $\ell$-integral with the same reduction as $\mathcal E(0)$.
So it suffices to show that $\mathcal E(0)$ is $\ell$-integral and $\det\overline{\mathcal E(0)}\ne0$.

\subsection{Step 3: splitting by progressions}
For $|v|<1$ and a polynomial $g$, $\sigma(v,g):=\sum_{j\ge0}v^jg(j)=\sum_b\Delta^bg(0)\,v^b/(1-v)^{b+1}$. Since $w$, $w^\ell$ and
$1-w^\ell\equiv1-w=(q^2-1)/q^2$ are $\ell$-units, we get the following. If $g$ has $\ell$-integral values on $\ZZ_{\ge0}$, then
$\sigma(w^\ell,g)$ is $\ell$-integral. If moreover $\bar g(j)=g_0(j)$ for all $j$, for a polynomial $g_0$ over $\FF_\ell$, then
$\overline{\sigma(w^\ell,g)}=\sum_b\Delta^bg_0(0)\,\bar w^b/(1-\bar w)^{b+1}$ (use $\bar w^\ell=\bar w$). Splitting $t\ge1$ into the
residue classes modulo $\ell$,
\[
  \mu(F)=\sum_{\gamma<\ell}w^{t_\gamma}\,\sigma\bigl(w^\ell,\ j\mapsto F'(t_\gamma+j\ell)\bigr),
\]
where the \emph{progression of $\gamma$} is $\{t\ge1:\ell\mid2t+2\gamma+1\}$ and $t_\gamma$ is its least element: $t_\gamma=\kappa-\gamma$
(type S), $\kappa-\gamma+\ell$ (type L), and $\ell$ (type D). On the progression of $\gamma$, $\ell\mid t-x_m$ iff $m\in C_\gamma$,
and $\ell\mid t$ iff $\gamma=\kappa$. Since $\ell^2>2h$, all these divisibilities are simple.

\subsection{Step 4: integrality and block triangularity}
We have $D'(x_i)=\prod_{m\ne i}(m-i)=\ell^2\,c_r^{-1}\,U_i$ with $U_i:=\prod_{m\notin C_\gamma}(x_i-x_m)$ an $\ell$-unit, so
$\nu_i=c_rU_i^{-1}$ is a unit. Let $\gamma_i,\gamma_j$ be the classes of $i,j$ and consider the contribution of the
progression of $\gamma$.

\emph{Off-diagonal entries.} At $t$ on the progression of $\gamma$, the factors of $F_{ij}$ divisible by $\ell$ number
$N_\gamma=3-\ind{\gamma=\gamma_i}-\ind{\gamma=\gamma_j}+\ind{\gamma=\kappa}$, and every term of $F'_{ij}(t)$ keeps at least $N_\gamma-1$ of
them. So the contribution has valuation $\ge N_\gamma-1-\delta_j\ge0$, with equality only if $\gamma=\gamma_i=\gamma_j$ (inside a
diagonal block) or $\gamma=\gamma_i\ne\gamma_j=\kappa$ (a non-D row against a D column).

\emph{Diagonal entries.} Off its own progression, $|t-x_i|_\ell=1$ and
$Q_i'(t)=\bigl(g_i'(t)(t-x_i)-g_i(t)+g_i(x_i)\bigr)/(t-x_i)^2$ with $g_i$ having at least three roots $\ell$-adically close to $t$,
so $v_\ell(Q_i'(t))\ge2$ and the contribution has valuation $\ge2-\delta_i\ge1$. The pole term
$4\ell^{2-\delta_i}x_iS_i(q)$ is integral, since $v_\ell(x_i)=\delta_i$ and the denominators $s^2$ of $S_i$ have $v_\ell(s)\le1$.
The own-progression contributions are integral by Step~5.

Hence $\mathcal E(0)$ is $\ell$-integral. Ordering the classes with $\kappa$ last, $\overline{\mathcal E(0)}$ is block upper triangular,
with $3\times3$ diagonal blocks $\overline{\mathcal E(0)}_{C_\gamma}$, and each diagonal block receives only its own progression
and the pole terms.

\subsection{Step 5: the diagonal blocks}
Fix a class $\gamma$ and $i=\gamma+r\ell$, $j=\gamma+s\ell$. On the progression of $\gamma$ we have $t\equiv x_i\equiv x_\gamma\pmod\ell$.

For $r\ne s$ let $c:=3-r-s$ be the rank of the third class member, and for $r=s$ let $c_1,c_2$ be the ranks of the
other two members; we write $x_{(c)}:=x_{\gamma+c\ell}$. For a polynomial $P$ let $Q[P]:=(P(u)-P(x_i))/(u-x_i)$, so that
$Q[PR]=Q[P]\,R+P(x_i)\,Q[R]$ and $Q_i=Q[g_i]$.

\emph{Types S and L.} If $r\ne s$, then
$F_{ij}=(u-x_{(c)})\cdot u\prod_{m\notin C_\gamma}(u-x_m)$, and modulo $\ell$, $F'_{ij}(t)\equiv t\prod_{m\notin C_\gamma}(t-x_m)\equiv x_iU_i$.
If $r=s$, write $g_i=P\cdot R$ with $P=(u-x_{(c_1)})(u-x_{(c_2)})$ and $R=u\prod_{m\notin C_\gamma}(u-x_m)$. Then
$Q_i'=Q[P]'R+Q[P]R'+P(x_i)Q[R]'$ with $Q[P]=u+x_i-x_{(c_1)}-x_{(c_2)}$. At $t$ we have $Q[P](t)\equiv0\pmod\ell$, and
$\ell^2\mid P(x_i)$, so $Q_i'(t)\equiv R(t)\equiv x_iU_i\pmod\ell$. In both cases $\nu_i(\cdot)\equiv c_r\bar x_\gamma$, a constant along
the progression. Summing with Step~3, and using $\bar w^\kappa=\bar q^{-(\ell-1)}=1$ and $\bar w^\ell=\bar w$, the progression
contributes $\bar w^{t_\gamma}c_r\bar x_\gamma/(1-\bar w)$, where $\bar w^{t_\gamma}=\bar w^{-\gamma}$ (S) or $\bar w^{1-\gamma}$ (L). For
the pole term only the odd $s=\ell s'$ survive in $\ell^2S_i$, and $q^\ell\equiv q$, so
$\overline{\ell^2S_i}=\bar q^{2i+1}\,T(M)$ with $M$ the largest odd number $\le(2i+1)/\ell$: $M=2r-1$ (S) or $2r+1$ (L). Since
$\bar q^{2i+1}=\bar q\,\bar w^{-\gamma}\bar w^{-r}$, the diagonal block is
\[
  \overline{\mathcal E(0)}_{C_\gamma}=\bar x_\gamma\,\bar w^{-\gamma}\,\overline{M_{\mathrm{type}(\gamma)}(q)} .
\]
\emph{Type D} ($\gamma=\kappa$). Here $x_i=-\ell(r+\frac12)$, the progression is $t=\ell(1+j')$, $j'\ge0$, and we put
$y=t/\ell+\frac12=j'+\frac32$; then $t=\ell(y-\frac12)$ and $t-x_{(c)}=\ell(y+c)$. Let
$R:=\prod_{m\notin C_\kappa}(u-x_m)$, so $R(t)\equiv R(x_i)=U_i\pmod\ell$. If $r\ne s$, two roots of $F_{ij}=u(u-x_{(c)})R$ are close
to $t$, and $F'_{ij}(t)=(2t-x_{(c)})R(t)+t(t-x_{(c)})R'(t)\equiv\ell(2y+c-\frac12)\,U_i\pmod{\ell^2}$. If $r=s$, write
$g_i=PR$ with $P=u(u-x_{(c_1)})(u-x_{(c_2)})$. Then $\ell^3\mid P(x_i)$; $Q[P](t)$ is a quadratic form in
$t,x_i,x_{(c_1)},x_{(c_2)}$, all divisible by $\ell$, so $Q[P](t)\equiv0\pmod{\ell^2}$; and $Q[P]'(t)=2t+x_i-x_{(c_1)}-x_{(c_2)}=\ell(2y+\frac52-2r)$ (using $c_1+c_2=3-r$).
Hence $Q_i'(t)\equiv\ell(2y+\frac52-2r)\,U_i\pmod{\ell^2}$. With $\nu_i=c_rU_i^{-1}$ and the column factor $\ell^{-1}$, the
residues are $c_r(2y+c-\frac12)=c_r(2j'+e_{rs})$ for $r\ne s$ and $c_r(2y+\frac52-2r)=c_r(2j'+e_{rr})$ for $r=s$. Summing
$c_r(2j'+e_{rs})$ against $\bar w^{1+j'}$ over $j'\ge0$ gives $\bar w\,c_r\bigl(e_{rs}/(1-\bar w)+2\bar w/(1-\bar w)^2\bigr)$.
The pole term $4\ell x_iS_i=-4(r+\frac12)\ell^2S_i$ reduces to
$-4(r+\frac12)\bar q^{2r+1}T(2r+1)$. So $\overline{\mathcal E(0)}_{C_\kappa}=\overline{M_D(q)}$.

\subsection{Step 6: conclusion}
By Steps 4 and 5,
\[
  \det\overline{\mathcal E(0)}=\prod_{\gamma\ne\kappa}\bigl(\bar x_\gamma\bar w^{-\gamma}\bigr)^3\,\overline{\det M_{\mathrm{type}(\gamma)}(q)}\cdot\overline{\det M_D(q)} .
\]
For $\gamma\ne\kappa$, $\bar x_\gamma=-(2\gamma+1)/2\ne0$, and the three determinants are units by hypothesis. So
$\det\overline{\mathcal E(a/b)}=\det\overline{\mathcal E(0)}\ne0$ in $\FF_\ell$, hence $\det\mathcal E(a/b)\ne0$ and
$\Delta_{3\ell}(a/b)\ne0$. This proves Theorem~\ref{thm:lemmaP}. \qed

\begin{remark}
The same analysis works for every $h$ with $h\bmod\ell\le\kappa$ and three or more atoms per class. One then needs a
fourth block type with four atoms, and an explicit prime-gap theorem to find suitable $\ell$ near $h/3$ for all large
$h$ [proof sketch]. For irrationality, infinitely many $h$ suffice, and $h=3\ell$ avoids both complications. With at
most two atoms per class the analogous block constants of type S vanish identically in our computations, so primes
$\ell>h/2$ did not give a usable nonvanishing statement [numerical].
\end{remark}

% ===========================================================================
\section{Proof of Theorem~\ref{thm:main}}\label{sec:proof}
% ===========================================================================

Let $q\in\{3,5\}$. We verify the hypotheses of Theorem~\ref{thm:conditional}.
\begin{enumerate}
\item Theorem~\ref{thm:arith} gives $d_h$ with $\alpha=\alpha_q$.
\item Theorem~\ref{thm:B}, applied to the certificate of Theorem~\ref{thm:cert}, gives (ii) with $\beta=-2.078$
  ($q=3$) or $\beta=-2.385$ ($q=5$).
\item Theorem~\ref{thm:nonvan} gives (iii).
\end{enumerate}
It remains to check $\alpha_q+\beta<0$, using only elementary bounds for logarithms: $0.6931471803<\log2<0.6931471808$,
and, from $\log x\ge1-1/x$, $\log3\ge\log2+\frac13$, $\log3\le2\log2-\frac14$ and $\log5\ge2\log2+\frac15$. Then
\[
  \alpha_3-2.078\le\frac{133}{60}+\frac{\log2}4-\frac{\log2}3-\frac19-2.078<-0.030,
\]
\[
  \alpha_5-2.385\le\frac{133}{60}+\frac{\log2}4+\frac{2\log2-\frac14}4-\frac{2\log2+\frac15}3-2.385<-0.23 .
\]
(The true values are $\alpha_3-2.078=-0.05425\ldots$ and $\alpha_5-2.385=-0.25687\ldots$.) Theorem~\ref{thm:conditional}
shows that $\chi_2(1/3)$ and $\chi_2(1/5)$ are irrational. \qed

% ===========================================================================
\section{The formalisation}\label{sec:lean}
% ===========================================================================

\subsection{What is formalised}
The repository \url{https://github.com/telaviv/chi2-irrationality} \cite{Chi2Repo} contains a proof of
Theorem~\ref{thm:main} in Lean~4 \cite{LeanFour} (version v4.35.0) with Mathlib \cite{Mathlib} (version v4.35.0
\cite{Mathlib4}). It has 61 Lean files under \lean{Chi2Irrational\slash{}}, with about $18{,}000$ lines. The formalisation covers every step of the proof,
including the parts that a written proof would cite: the prime number theorem, the Andreief identity, the lattice
energy inequality, both dual certificates and the block determinants. Paths below are relative to the folder
\lean{Chi2Irrational\slash{}}, and all declarations are in the namespace \lean{Chi2Irr}.

\subsection{The statement}
The definitions that the reader has to check are, in \lean{Defs.lean} and \lean{Statement.lean}:
\begin{quote}\ttfamily\footnotesize
noncomputable def chi2Term (z : $\mathbb{R}$) (n : $\mathbb{N}$) : $\mathbb{R}$ :=\\
\hspace*{2em}z \^{} (2 * n + 1) / (2 * (n : $\mathbb{R}$) + 1) \^{} 2\\[2pt]
noncomputable def chi2 (z : $\mathbb{R}$) : $\mathbb{R}$ :=\\
\hspace*{2em}$\sum$' n : $\mathbb{N}$, chi2Term z n\\[2pt]
def Chi2OneThirdIrrational : Prop := Irrational (chi2 (1 / 3))\\
def Chi2OneFifthIrrational : Prop := Irrational (chi2 (1 / 5))\\[2pt]
theorem chi2\_one\_third\_irrational : Chi2OneThirdIrrational\\
theorem chi2\_one\_fifth\_irrational : Chi2OneFifthIrrational
\end{quote}
The theorems are in \lean{Arith\slash{}Assembly.lean}. Here $\sum'$ is Mathlib's unconditional sum (\lean{tsum}).
\lean{Defs.lean} proves that the series is summable for $|z|\le1$ (\lean{summable\_chi2Term}), so \lean{chi2 z} is the
genuine value $\chi_2(z)$. \lean{Irrational x} is Mathlib's predicate $x\notin\QQ$.

The command \lean{lake env lean Axioms.lean} prints
\begin{quote}\ttfamily\small
'Chi2Irr.chi2\_one\_third\_irrational' depends on axioms:\\
\hspace*{2em}[propext, Classical.choice, Quot.sound]\\
'Chi2Irr.chi2\_one\_fifth\_irrational' depends on axioms:\\
\hspace*{2em}[propext, Classical.choice, Quot.sound]
\end{quote}
and the same list for each of the main ingredients listed in \lean{Axioms.lean}. These are Lean's three standard axioms,
which Mathlib also uses. The development contains no \lean{sorry}, no added axiom, no \lean{native\_decide} and no
\lean{implemented\_by}, so no step relies on compiled code.

\subsection{Correspondence with the paper}
Table~\ref{tab:lean} lists the main formal counterparts.

\begin{table}[ht]
\centering\small
\begin{tabular}{@{}>{\raggedright\arraybackslash}p{0.27\textwidth}>{\raggedright\arraybackslash}p{0.69\textwidth}@{}}
\toprule
paper & Lean declaration (file) \\
\midrule
Lemma~\ref{lem:pole}, \eqref{eq:value} & \lean{powSum\_eq\_tsum}, \lean{pole\_series} (\lean{PoleSeries}); \lean{Lam\_eq}, \lean{aeval\_chi2\_Delta} (\lean{PoleExpansion})\\
Prop.~\ref{prop:criterion}, Thm.~\ref{thm:conditional} & \lean{irrational\_of\_poly\_approx} (\lean{Criterion}); \lean{MachineBounds}, \lean{irrational\_of\_bounds} (\lean{MainCond})\\
Section~\ref{sec:free} & \lean{map\_DeltaF} (\lean{Arith\slash{}Free}); \lean{det\_GLeX}, \lean{det\_GramF} (\lean{Arith\slash{}Gauss})\\
Lemma~\ref{lem:nine} & \lean{lemma9\_DeltaF} (\lean{Arith\slash{}Lemma9})\\
Prop.~\ref{prop:classbound}, Lemma~\ref{lem:classcount} & \lean{classBound\_DeltaF} (\lean{Arith\slash{}ClassBound}); \lean{classBd\_le\_phi1}, \lean{classBd\_eq\_zero} (\lean{Arith\slash{}ClassCount})\\
Thm.~\ref{thm:generic} & \lean{eventually\_primeSum\_phi1\_le}, \lean{eventually\_primeSum\_lev\_le} (\lean{Arith\slash{}PrimeSums}); \lean{PNT.theta\_div\_tendsto} (\lean{PNT\slash{}Theta})\\
Lemma~\ref{lem:poleorder}, \eqref{eq:clear} & \lean{EDeg\_DeltaF}, \lean{aeval\_clearE} (\lean{Arith\slash{}PoleOrder})\\
Lemma~\ref{lem:qadic}, \eqref{eq:andreief} & \lean{ord\_DeltaP}, \lean{qdeg\_clearE\_DeltaF} (\lean{Arith\slash{}QAdic}); \lean{factorial\_smul\_det\_sum\_rankOne} (\lean{Analytic\slash{}Andreief})\\
Lemmas~\ref{lem:two}, \ref{lem:fdseq} & \lean{twoAdic\_Delta} (\lean{Arith\slash{}TwoAdic}); \lean{V2\_fd\_seq}, \lean{V2\_U} (\lean{Arith\slash{}TwoAdicBasis}); \lean{V2\_entry} (\lean{Arith\slash{}TwoAdicEntries})\\
Thm.~\ref{thm:arith} & \lean{denQ}, \lean{padicNorm\_Delta\_eN}, \lean{denQ\_int} (\lean{Arith\slash{}Special}); \lean{arith\_denQ} (\lean{Arith\slash{}Assembly})\\
node expansion, Lemma~\ref{lem:slotdet} & \lean{HT\_rankOne}, \lean{abs\_Delta\_le\_of} (\lean{Analytic\slash{}Nodes}); \lean{abs\_det\_slot\_le} (\lean{Analytic\slash{}SlotDet})\\
Lemmas~\ref{lem:polya}, \ref{lem:dual} & \lean{polya\_log} (\lean{Analytic\slash{}Polya}); \lean{pattern\_le}, \lean{abs\_det\_HT\_le\_dual} (\lean{Analytic\slash{}Dual})\\
Lemma~\ref{lem:discrete} & \lean{log\_bN\_le}, \lean{StepData.Pm\_le}, \lean{StepData.Em\_ge} (\lean{Analytic\slash{}Discrete})\\
Def.~\ref{def:cert}, Thm.~\ref{thm:B} & \lean{DualCert}, \lean{DualCert.det\_bound}, \lean{DualCert.analytic} (\lean{Analytic\slash{}TheoremB})\\
Thm.~\ref{thm:cert} & \lean{CData.Valid}, \lean{CData.toCert}, \lean{CData.Phi\_le} (\lean{Cert\slash{}Instance}); \lean{Cert.Q3.analytic}, \lean{Cert.Q5.analytic} (\lean{Cert\slash{}Q3}, \lean{Cert\slash{}Q5})\\
Thm.~\ref{thm:lemmaP} & \lean{Delta\_eval\_ne\_zero\_of\_det} (\lean{NonVan\slash{}Gram}); \lean{MS}, \lean{ML}, \lean{MD} (\lean{NonVan\slash{}BlockMat}); \lean{lemmaP} (\lean{NonVan\slash{}LemmaP})\\
Thm.~\ref{thm:nonvan} & \lean{nonvan\_of}, \lean{nonvan\_three}, \lean{nonvan\_five} (\lean{NonVan\slash{}NonVanQ})\\
Section~\ref{sec:proof} & \lean{margin\_three}, \lean{margin\_five}, \lean{boundsThree}, \lean{boundsFive}, \lean{chi2\_one\_third\_irrational}, \lean{chi2\_one\_fifth\_irrational} (\lean{Arith\slash{}Assembly})\\
\bottomrule
\end{tabular}
\caption{Formal counterparts of the main steps.}\label{tab:lean}
\end{table}

A few arguments are organised differently in the formalisation.
\begin{itemize}
\item The prime number theorem $\theta(x)/x\to1$ is ported from the PrimeNumberTheoremAnd project \cite{PNTplus}
  (Wiener--Ikehara route), through the port in the \lean{zeta5} repository of C\'esar Del Solar \cite{zeta5} and the
  author's formalisation \cite{BetaRepo} of a result on Dirichlet beta values. The ported files keep their Apache
  License~2.0.
\item The discretization bounds of Lemma~\ref{lem:discrete}(ii)--(iii) are proved piece by piece of the step density,
  with the antiderivatives $\psi$ and $\Psi$, and with the error terms $8J(1+\log h)+2\log(Nh+1)$ and $4h\log(Nh+1)$
  instead of those stated here. Only their size $O(h\log h)$ matters.
\item The value $\beta_h=\lfloor h^2/4\rfloor$ of Lemma~\ref{lem:poleorder} is not formalised; the formal proof uses
  $4\beta_h\le h^2+2h$.
\item The bound $a_h\ge(h^2-7h-2)/3$ and the two margins are proved with the elementary logarithm bounds of
  Section~\ref{sec:proof}. The formal $a_h$ is $(h(h-1)-6h)/3$ in natural-number arithmetic, which is the
  $\max(0,\cdot)$ of Section~\ref{sec:arith2}.
\item The potential $U^\mu$ and the energy $E(\mu)$ of a step density are \emph{defined} in Lean by their closed forms
  $\sum_jc_j(\psi(x-a_j)-\psi(x-a_{j+1}))$ and $\sum_{j,j'}c_jc_{j'}(\cdots)$ with $\Psi$, which equal the integrals of
  Section~\ref{sec:analytic}. Both Theorem~\ref{thm:B} and the certificates use these closed forms.
\end{itemize}

\subsection{How the certificates are checked}
The dual certificates are records of rational data (\lean{Cert\slash{}DataQ3.lean}, \lean{Cert\slash{}DataQ5.lean}): the
log table as a list of integers in units of $2^{-40}$, the heights $c_j$, the bounds $S_k$, $\lambda$ and $E_{\mathrm{lo}}$.
\lean{CData.Valid} collects every check of Section~\ref{sec:cert} as a Boolean computation, and \lean{Cert\slash{}Q3.lean}
and \lean{Cert\slash{}Q5.lean} prove it by \lean{decide +kernel}, i.e.\ by evaluation in the Lean kernel. The interval
arithmetic (\lean{Interval\slash{}}) has exact rational endpoints (only the log-table entries and the enclosure of $\log2$
are dyadic by construction), and its soundness lemmas turn
\lean{check = true} into real inequalities; the enclosure of $\log2$ is itself checked by \lean{decide +kernel}
against Mathlib's series bound. The data were produced by the script \lean{scripts\slash{}gen\_cert.py}, which is not
trusted: a wrong value would make the build fail but could not make a false statement provable. Re-running the
script reproduces both data files byte for byte. The six block determinants of Theorem~\ref{thm:nonvan} are evaluated
exactly by \lean{norm\_num}.

\subsection{Cost}
With Mathlib's prebuilt cache, a clean build of the project takes about ten minutes on a desktop computer. The two
kernel-checked certificates are a small part of it.

\subsection{What has to be trusted}
Beyond the Lean kernel and the three axioms, the reader has to trust that the definitions displayed above express
Theorem~\ref{thm:main}, and the Mathlib definitions they use (real numbers, \lean{tsum}, \lean{Irrational}). Nothing
else, in particular no numerical data and no generating script, needs to be trusted. To avoid trusting the
elaborator's handling of the environment, we replayed every declaration of the $61$ project modules through the
kernel with \lean{leanchecker} (the \lean{lean4checker} tool shipped with the Lean toolchain), run on each module
against the compiled environment of its imports; it reported no errors. Mathlib itself was not replayed from scratch
(\lean{--fresh}). Apart from the check
by the Lean kernel, all verification of the formalisation so far was carried out by AI systems; no human expert has
yet reviewed the Lean code.

% ===========================================================================
\section{Remarks}\label{sec:remarks}
% ===========================================================================

\subsection{The margins}
For $q=3$ the exponent that must be negative is $\alpha_3+\Phi_{\mathbb Q}\approx-0.0546$ per $h^2$ with the certificate
value $\Phi_{\mathbb Q}$ ($-0.0543$ with the rounded value $-2.078$ used in Section~\ref{sec:proof}). The arithmetic part is
$2.02375$, so the margin is under $3\%$ of it. A $2$-adic constant larger than about $0.33$ instead of $\frac14$ in
Lemma~\ref{lem:two} would have consumed it. For $q=5$ the margin is $-0.257$, and no single constant is critical. With the
sharpest available analytic constant, the equilibrium value $C(w)$, the margins would be about $-0.077$ and $-0.328$
[numerical].

\subsection{Effectivity}
Every bound in the proof is effective except the prime number theorem, which is used without an error term in
Theorem~\ref{thm:generic}. Even with an explicit form, the $O(h\log h)$ terms and the prime-power contributions of
Lemma~\ref{lem:nine} decay slowly relative to the margin for $q=3$. The irrationality statement is not affected.

\subsection{Other values of \texorpdfstring{$q$}{q}}
The construction works verbatim for every integer $q\ge2$. The arithmetic part is
$\frac{133}{60}+c_2(q)\log2+\sum_{p\mid q^2-1,\,p\text{ odd}}\frac{v_p(q^2-1)}4\log p-\frac13\log q$, where $c_2(q)$ is the
$2$-adic constant ($c_2=\frac14$ when $v_2(q^2-1)=3$). The analytic part improves as $q$ grows [numerical]. We expect the method to prove $\chi_2(1/q)\notin\QQ$ for
every $q\ge3$ once these constants are made explicit, with $q=3$ the hardest case [heuristic]. We have not carried
this out.

\subsection{\texorpdfstring{$\chi_2(1/2)$}{chi2(1/2)}}
At $q=2$ ($w=\frac14$) the same construction fails [numerical]. The analytic part is only about $C(\frac14)\approx-1.744$,
the generic primes still cost $\frac{133}{60}$, and the prime $3\mid q^2-1$ costs $\frac14\log3$. The $2$-adic gain
($-\frac13\log2$) does not compensate. The limiting exponent is about $+0.52$ per $h^2$, and exact computations up to
$h=64$ confirm the trend. Variations of the external field (pole sets, numerator zeros, node regions) gained at most
$0.01$ in our experiments. Since $\chi_2(\frac13)+\chi_2(\frac12)=\frac{\pi^2}8-\frac12\log2\log3$, a proof that
$\chi_2(1/2)$ is irrational, together with Theorem~\ref{thm:main}, would still not decide the irrationality of
$\log2\log3$.

\section*{Acknowledgements and disclosure}
The construction, the proofs, the certificate programs, the Lean formalisation and the text of this paper were produced
largely by AI models (Claude, Anthropic), working under the direction of the author. Apart from the check by the Lean
kernel, all verification so far, including the internal reviews of the argument, was carried out by AI systems. No
human expert has yet reviewed the mathematics or the Lean code; we would be grateful for corrections. The prime number
theorem in the formalisation is ported from the PrimeNumberTheoremAnd project of A.~Kontorovich, T.~Tao and
collaborators, through the \lean{zeta5} repository of C\'esar Del Solar.

\bibliographystyle{amsplain}
\bibliography{references}

\end{document}
